% ============================================================================
% GUÍA DE USO — Análisis de rentabilidad del negocio de paquetería
%
% Archivo principal. Cada capítulo está en la carpeta capitulos/.
% Para compilar:  latexmk -pdf guia_usuario.tex   (o: build/construir_guia.sh)
%
% · Los números de los ejercicios salen de valores.tex y las imágenes de img/.
%   Ambos los genera build/capturas.py a partir del Excel: si el Excel cambia,
%   vuelva a ejecutar ese script y recompile; la guía se actualiza sola.
% · Los colores y cuadros están en estilo.sty.
% ============================================================================
\documentclass[11pt,oneside]{report}
\usepackage{estilo}
% Archivo generado por build/capturas.py. No editar a mano.

\newcommand{\vUAItotal}{\$9,168}
\newcommand{\vUNtotal}{\$5,959}
\newcommand{\vIngtotal}{\$32,230}
\newcommand{\vMCpct}{59.6\%}
\newcommand{\vCostoKg}{\$0.575}
\newcommand{\vIngKg}{\$0.804}
\newcommand{\vPEkg}{16,606}
\newcommand{\vMSeg}{58.6\%}
\newcommand{\vUAIocc}{\$8,420}
\newcommand{\vUAIcen}{\$862}
\newcommand{\vUAIori}{(\$114)}
\newcommand{\vUAIcUno}{\$2,175.82}
\newcommand{\vUAIcTres}{\$3,043.80}
\newcommand{\vAlertaTres}{Combustible distribución Oriente +44\% sobre norma.}
\newcommand{\vUAIsep}{\$4,070}
\newcommand{\vUAIoct}{\$5,099}
\newcommand{\vFijosMes}{\$3,240.07}
\newcommand{\vPEsep}{8,643}
\newcommand{\vPEcontSep}{0.9}
\newcommand{\vSimUAIcont}{\$2,376}
\newcommand{\vSimTminOri}{\$0.796}
\newcommand{\vSimTminOcc}{\$0.281}
\newcommand{\vSimUAImes}{\$4,752}
\newcommand{\vSimPE}{0.81}
\newcommand{\vSimGrid}{\$4,752}
\newcommand{\vSimGridMas}{\$3,537}
\newcommand{\vAlertasIni}{1}
\newcommand{\vAbiertosIni}{1}
\newcommand{\vSimUAIori}{\$9}
\newcommand{\vSimMgOri}{0.4\%}
\newcommand{\vEjDosUAI}{\$2,969.24}
\newcommand{\vEjDosUAItres}{\$3,632.01}
\newcommand{\vEjTresUAI}{\$2,806.24}
\newcommand{\vEjTresAlerta}{Combustible traslado Oriente +16\% sobre norma.}
\newcommand{\vEjTresAlertas}{2}
\newcommand{\vEjCuatroUAIcinco}{\$2,806.24}
\newcommand{\vEjCuatroUAIseis}{(\$866.67)}
\newcommand{\vEjCuatroCombSeis}{\$138.00}
\newcommand{\vEjCuatroCombCinco}{\$128.00}
\newcommand{\vEjCincoDif}{(\$8,000.00)}
\newcommand{\vEjCincoReal}{\$24,480.00}
\newcommand{\vEjCincoFijosOct}{\$3,255.07}
\newcommand{\vEjCincoUAIoct}{\$8,949}
\newcommand{\vEjCincoUAInov}{(\$867)}
\newcommand{\vEjCincoPagoTotal}{\$1,784.40}
\newcommand{\vEjCincoPagoTotalAntes}{\$1,544.40}
\newcommand{\vEjCincoDeclarada}{\$16,480.00}
\newcommand{\vEjCincoHabAntes}{3334.74}
\newcommand{\vEjCincoHab}{7334.74}
\newcommand{\vEjCincoVcAntes}{1005.71}
\newcommand{\vEjCincoVc}{3005.71}
\newcommand{\vEjCincoScAntes}{1575.00}
\newcommand{\vEjCincoSc}{3575.00}
\newcommand{\vEjCincoFijosOctAntes}{\$3,240.07}
\newcommand{\vEjOchoUAI}{(\$3,240)}
\newcommand{\vEjOchoConc}{CUADRA}
\newcommand{\vEjOchoFijos}{\$3,240.07}
\newcommand{\vEjSeisCont}{4}
\newcommand{\vEjSeisUAI}{\$8,082}
\newcommand{\vEjSeisMargen}{16.9\%}
\newcommand{\vEjSieteUAIori}{\$367}
\newcommand{\vEjSieteMgOri}{13.7\%}
\newcommand{\vEjSieteUAImes}{\$5,468}


\newcommand{\libro}{\textsf{Analisis\_\allowbreak Rentabilidad\_\allowbreak Paqueteria.xlsx}}
\newcommand{\version}{1.0 · octubre de 2026}

\hypersetup{pdftitle={Guía de uso – Análisis de rentabilidad de paquetería},
            pdfsubject={Manual paso a paso del libro de Excel}}

\begin{document}
% ---------------------------------------------------------------- PORTADA
\begin{titlepage}
\begin{tikzpicture}[remember picture,overlay]
  \fill[azul] (current page.north west) rectangle ([yshift=-9cm]current page.north east);
  \fill[naranja] ([yshift=-9cm]current page.north west) rectangle ([yshift=-9.25cm]current page.north east);
  \node[anchor=north west,text=white,font=\sffamily\bfseries\fontsize{30}{36}\selectfont,align=left]
    at ([xshift=2.2cm,yshift=-2.6cm]current page.north west) {Guía de uso};
  \node[anchor=north west,text=white,font=\sffamily\Large,align=left,text width=16cm]
    at ([xshift=2.2cm,yshift=-4.3cm]current page.north west)
    {Libro de Excel para el\\[2pt]\textbf{Análisis de rentabilidad del negocio de paquetería}};
  \node[anchor=north west,text=white!85!azul,font=\sffamily\normalsize]
    at ([xshift=2.2cm,yshift=-7.2cm]current page.north west) {Manual paso a paso, con imágenes y ejercicios de práctica};
\end{tikzpicture}

\vspace*{8.2cm}
{\large\color{azul}\textbf{Para quién es esta guía}}\par\smallskip
Para la persona que llevará el libro de Excel. No hace falta saber Excel ni contabilidad: todo se explica desde el principio.
Si sigue los capítulos en orden y hace los ejercicios, en una o dos sesiones de trabajo podrá registrar contenedores,
hacer el cierre de cada mes e interpretar los resultados.

\bigskip
{\large\color{azul}\textbf{Cómo leer esta guía}}\par\smallskip
\begin{tabularx}{0.96\linewidth}{@{}lX@{}}
\recuadro{1} & Los círculos rojos numerados de las imágenes indican, en orden, dónde debe mirar o escribir.\\[3pt]
\hoja{ENTRADA} & Nombre de una hoja (pestaña) del libro.\\[3pt]
\celda{C5} & Dirección de una celda: columna C, fila 5.\\[3pt]
\escriba{5000} & Lo que debe escribir, tal cual.\\[3pt]
\tecla{Enter} & Una tecla del teclado.\\
\end{tabularx}

\medskip
\begin{tcolorbox}[colback=green!8,colframe=verde,boxrule=0.6pt,arc=2pt,left=6pt,right=6pt,top=3pt,bottom=3pt]
\small Los cuadros verdes \textbf{«Lo que debe ver si lo hizo bien»} traen los números exactos que deben aparecer en su pantalla
después de cada ejercicio. Si sus números coinciden, lo hizo bien.
\end{tcolorbox}

\vfill
{\small\color{gris}
Archivo del libro: \libro\hfill Versión de la guía: \version\\
Las imágenes y los números de esta guía se generan automáticamente a partir del propio libro de Excel.}
\end{titlepage}

\pagenumbering{roman}
\tableofcontents
\clearpage\pagenumbering{arabic}
% ============================================================================
\chapter{Antes de empezar}
% ============================================================================

\section{¿Qué hace este libro de Excel?}

El libro responde, de forma automática, a tres preguntas:

\begin{enumerate}
  \item \textbf{¿Cuánto ganamos o perdemos con cada contenedor?}
  \item \textbf{¿Qué región deja dinero y cuál no?} (Occidente, Centro y Oriente).
  \item \textbf{¿Cómo va cada mes?}: si se cubren los gastos fijos, cuántos kilogramos hacen falta para no perder (punto de equilibrio)
        y cuánto se pagará de impuestos.
\end{enumerate}

Usted solo escribe los datos de cada contenedor y algunos datos del mes. El libro hace todas las cuentas: más de
32\,000 fórmulas que ya están escritas y protegidas.

\section{El negocio, tal como lo calcula el libro}

El libro sigue el mismo recorrido que hace la mercancía. Cada etapa tiene sus gastos:

\begin{figure}[H]\centering
\begin{tikzpicture}[
  etapa/.style={draw=azul,fill=azulclaro,rounded corners=3pt,text width=3.25cm,align=center,minimum height=2.6cm,font=\small},
  flecha/.style={-{Stealth[length=3mm]},very thick,azul}]
\node[etapa] (p) {\textbf{1. PUERTO}\\[2pt]Vehículo propio\\\footnotesize cita (10 USD), 60 L de combustible, dieta del chofer, desgaste};
\node[etapa,right=0.55cm of p] (t) {\textbf{2. TRANSITARIA}\\[2pt]ENCI o A.V\\\footnotesize desagrupe pagado por jornal, estadía};
\node[etapa,right=0.55cm of t] (d) {\textbf{3. DISTRIBUCIÓN}\\[2pt]Servicio de terceros\\\footnotesize chofer, ayudante, combustible; en Centro y Oriente también traslado y 2.º desagrupe};
\node[etapa,right=0.55cm of d] (c) {\textbf{4. COMERCIALIZACIÓN}\\[2pt]21 gestores\\\footnotesize pago fijo igual para todos + parte variable};
\draw[flecha] (p) -- (t); \draw[flecha] (t) -- (d); \draw[flecha] (d) -- (c);
\node[below=0.35cm of t.south east,anchor=north,text width=14cm,align=center,font=\small\color{gris}]
 {Además: gastos fijos del mes (oficina, salarios, alquiler, depreciación del camión…) y tributos (10\,\% ventas y servicios,
 1\,\% territorial, 14\,\% seguridad social, 5\,\% fuerza de trabajo y 35\,\% utilidades).};
\end{tikzpicture}
\caption{Recorrido de la mercancía y gastos de cada etapa}
\end{figure}

Con los ingresos (kilogramos $\times$ tarifa por kg de cada región) y todos esos gastos, el libro calcula la
\textbf{utilidad} (ganancia) de cada contenedor, de cada región y de cada mes.

\section{Qué necesita}

\begin{itemize}
  \item Una computadora con \textbf{Microsoft Excel 2010 o más reciente}. También funciona con LibreOffice Calc (gratuito).
  \item El archivo \libro. No tiene macros, así que no hace falta «habilitar contenido».
  \item Los documentos de cada contenedor: manifiesto o desglose por destino, vales de combustible, hojas de ruta,
        comprobantes de pago de la cita, de la dieta y de los jornales.
  \item Para el cierre del mes: facturas y nóminas de los gastos fijos y lo cobrado por cada gestor.
\end{itemize}

\section{Su papel: una sola persona lleva el libro}

Todo el libro lo lleva \textbf{una sola persona}: usted. Eso simplifica las cosas, pero obliga a tener orden.
Su trabajo se reduce a cuatro momentos, que la propia hoja \hoja{INICIO} le recuerda:

\begin{figure}[H]\centering
\begin{tikzpicture}[
  caja/.style={draw=azul,rounded corners=3pt,text width=6.6cm,minimum height=1.55cm,align=left,font=\small,inner sep=5pt},
  tit/.style={font=\small\bfseries\color{white},fill=azul,rounded corners=2pt,inner sep=3pt}]
\node[caja] (a) {\textbf{Una sola vez}, al empezar.\\Cargar sus datos reales y borrar los ejemplos. \hfill\textit{Capítulo \ref{cap:preparar}}};
\node[caja,right=0.5cm of a] (b) {\textbf{Con cada contenedor}, en dos momentos:\\al llegar y al terminar la distribución. \hfill\textit{Capítulo \ref{cap:diario}}};
\node[caja,below=0.45cm of a] (c) {\textbf{Cada fin de mes}: revisar, completar gastos reales, liquidar a los gestores y guardar una copia. \hfill\textit{Capítulo \ref{cap:cierre}}};
\node[caja,below=0.45cm of b] (d) {\textbf{Cuando algo cambia}: un precio, un vehículo, un gestor, un impuesto. \hfill\textit{Capítulo \ref{cap:cambios}}};
\end{tikzpicture}
\caption{Los cuatro momentos de trabajo con el libro}
\end{figure}

\begin{consejo}
Como nadie más revisa sus anotaciones, el libro trae \textbf{controles automáticos} que hacen de «segunda persona».
Avisan si falta un dato, si el combustible se pasa de la norma o si algo no cuadra. Consúltelos siempre en la hoja
\hoja{INICIO} antes de dar un mes por cerrado.
\end{consejo}

\section{Ruta de aprendizaje recomendada}

\begin{enumerate}
  \item Si nunca ha usado Excel, lea el capítulo~\ref{cap:excel} (unos 20 minutos).
  \item Recorra el libro con el capítulo~\ref{cap:conozca}, que explica qué hace cada hoja, con imágenes.
  \item Haga \textbf{los 8 ejercicios del capítulo~\ref{cap:taller}} sobre una copia de práctica. Es la parte más importante:
        al terminar habrá hecho, una vez, todo lo que hará en su trabajo real.
  \item Prepare el libro con sus datos reales (capítulo~\ref{cap:preparar}) y empiece a trabajar.
  \item Los capítulos \ref{cap:diario} a \ref{cap:problemas} y los anexos son material de consulta para el día a día.
\end{enumerate}

\begin{atencion}
\textbf{Los datos que trae el libro son de EJEMPLO}: 4 contenedores de septiembre y octubre de 2026, una tarifa de 0,80 USD/kg y
salarios y normas inventados. Sirven para aprender y para los ejercicios. Antes de usar el libro de verdad hay que sustituirlos
siguiendo el capítulo~\ref{cap:preparar}. Solo son datos reales de la empresa la cita de 10 USD y los 60 litros por viaje al puerto.
\end{atencion}

% ============================================================================
\chapter{Excel desde cero: lo mínimo que necesita saber}\label{cap:excel}
% ============================================================================

Si ya usa Excel con soltura, puede pasar al capítulo~\ref{cap:conozca}. Si no, este capítulo le explica, sin tecnicismos,
todo lo que hace falta para trabajar con el libro.

\section{Las partes de la pantalla}

\begin{figure}[H]\centering
\begin{tikzpicture}[x=1cm,y=1cm,font=\footnotesize]
% ventana
\draw[gris,fill=white] (0,0) rectangle (15,9.2);
% barra de título y cinta
\fill[verde!70!black] (0,8.6) rectangle (15,9.2);
\node[white,anchor=west] at (0.2,8.9) {\bfseries Analisis\_Rentabilidad\_Paqueteria.xlsx — Excel};
\fill[grisclaro] (0,7.3) rectangle (15,8.6);
\foreach \x/\t in {0.3/Archivo,1.5/Inicio,2.6/Insertar,3.8/Diseño,4.9/Fórmulas,6.1/Datos,7.1/Revisar,8.3/Vista}
  \node[anchor=west] at (\x,8.35) {\t};
\foreach \x in {0.4,1.5,2.6,3.7,4.8,5.9,7.0,8.1,9.2}
  \draw[gris!60,fill=white] (\x,7.45) rectangle ++(0.9,0.65);
% barra de fórmulas
\draw[gris,fill=white] (0,6.75) rectangle (15,7.25);
\draw[gris] (1.6,6.75) -- (1.6,7.25);
\node[anchor=west] at (0.1,7.0) {\bfseries C8};
\node[anchor=west] at (1.9,7.0) {06/10/2026};
% cuadrícula
\fill[grisclaro] (0,1.35) rectangle (0.6,6.75);
\fill[grisclaro] (0.6,6.35) rectangle (15,6.75);
\foreach \c [count=\i] in {A,B,C,D,E,F,G,H}
  \node at (0.6+\i*1.6-0.8,6.55) {\c};
\foreach \r [count=\i] in {1,...,10}
  \node at (0.3,6.35-\i*0.5+0.25) {\r};
\foreach \i in {0,...,8} \draw[gris!40] (0.6+\i*1.6,1.35) -- (0.6+\i*1.6,6.75);
\foreach \i in {0,...,10} \draw[gris!40] (0,6.35-\i*0.5) -- (15,6.35-\i*0.5);
% celda activa C8
\draw[verde!70!black,line width=1.6pt] (3.8,2.35) rectangle (5.4,2.85);
\node at (4.6,2.6) {06/10/2026};
% triángulo de comentario
\fill[rojo] (2.2,6.35) -- (2.05,6.35) -- (2.2,6.2) -- cycle;
% desplegable
\draw[gris,fill=white] (5.42,2.36) rectangle (5.75,2.84); \node at (5.585,2.6) {\tiny$\blacktriangledown$};
% pestañas
\fill[grisclaro] (0,0.75) rectangle (15,1.35);
\foreach \x/\t in {0.9/INICIO,2.5/GUIA\_USO,4.3/ANALISIS,6.2/DASHBOARD}
  \node[draw=gris!60,fill=white,inner sep=2pt] at (\x,1.05) {\t};
\node[draw=azul,fill=white,inner sep=2pt,font=\footnotesize\bfseries] at (8.1,1.05) {ENTRADA};
\node at (9.4,1.05) {\dots};
\fill[grisclaro!70] (0,0) rectangle (15,0.75);
\node[anchor=east] at (14.8,0.37) {zoom $-$ \rule{1.5cm}{0.4pt} $+$ 90\%};
% rótulos
\tikzset{rot/.style={circle,fill=rojo,text=white,inner sep=1pt,minimum size=1.3em,font=\scriptsize\bfseries}}
\node[rot] at (11.6,8.9) {1};
\node[rot] at (10.6,7.8) {2};
\node[rot] at (1.2,6.95) {3};
\node[rot] at (9.6,6.55) {4};
\node[rot] at (0.3,1.85) {5};
\node[rot] at (5.9,2.95) {6};
\node[rot] at (5.95,2.6) {7};
\node[rot] at (2.4,6.1) {8};
\node[rot] at (11.2,1.05) {9};
\node[rot] at (12.6,0.37) {10};
\end{tikzpicture}
\caption{Esquema de la ventana de Excel (los colores varían según la versión)}
\end{figure}

\begin{tabularx}{\linewidth}{@{}c>{\bfseries}l X@{}}
\recuadro{1} & Barra de título & Muestra el nombre del archivo abierto. \\
\recuadro{2} & Cinta de opciones & Los menús: Archivo, Inicio, Revisar\dots{} Con este libro casi no hará falta usarlos. \\
\recuadro{3} & Cuadro de nombres & Indica en qué celda está. Aquí, \celda{C8}. \\
\recuadro{4} & Letras de columna & Las columnas se nombran con letras: A, B, C\dots{} Después de Z siguen AA, AB, AC\dots \\
\recuadro{5} & Números de fila & Las filas se numeran: 1, 2, 3\dots \\
\recuadro{6} & Celda activa & La celda seleccionada, con un borde grueso. Lo que escriba irá ahí. \\
\recuadro{7} & Flecha de lista & Aparece en algunas celdas: al pulsarla muestra las opciones permitidas. \\
\recuadro{8} & Triángulo rojo & La celda tiene una nota de ayuda. Ponga el ratón encima para leerla. \\
\recuadro{9} & Pestañas de hoja & Cada pestaña es una hoja del libro. Haga clic para cambiar de hoja. \\
\recuadro{10} & Zoom & Agranda o achica la vista. No cambia los datos. \\
\end{tabularx}

\section{Libro, hoja, fila, columna y celda}

\begin{itemize}
  \item \textbf{Libro}: el archivo completo (\libro).
  \item \textbf{Hoja}: cada «página» del libro. Este libro tiene 14, por ejemplo \hoja{INICIO}, \hoja{ENTRADA} o \hoja{PRECIOS}.
  \item \textbf{Celda}: cada casilla de la cuadrícula. Se nombra con la letra de su columna y el número de su fila.
        \celda{B10} es la celda de la columna B y la fila 10.
\end{itemize}

\begin{ejemplo}
Cuando la guía dice «escriba \escriba{5000} en la celda \celda{G10} de la hoja \hoja{ENTRADA}», significa:
\begin{pasos}
\paso Haga clic en la pestaña \hoja{ENTRADA}, abajo.
\paso Busque la columna \textbf{G} (arriba) y la fila \textbf{10} (a la izquierda). Haga clic en la casilla donde se cruzan.
\paso Escriba \escriba{5000} y pulse \tecla{Enter}.
\end{pasos}
\end{ejemplo}

\section{Cómo moverse}

\begin{tabularx}{\linewidth}{@{}l X@{}}
\toprule
\textbf{Para\dots} & \textbf{Haga esto} \\ \midrule
Ir a una celda & Haga clic en ella. O pulse \tecla{F5}, escriba la dirección (por ejemplo \escriba{AB10}) y pulse \tecla{Enter}. \\
Moverse una celda & Use las flechas \tecla{$\leftarrow$} \tecla{$\rightarrow$} \tecla{$\uparrow$} \tecla{$\downarrow$}. \\
Pasar a la celda de la derecha después de escribir & \tecla{Tab} \\
Pasar a la celda de abajo después de escribir & \tecla{Enter} \\
Ir al principio de la hoja & \tecla{Ctrl}+\tecla{Inicio} \\
Cambiar a la hoja siguiente / anterior & \tecla{Ctrl}+\tecla{AvPág} / \tecla{Ctrl}+\tecla{RePág}, o clic en la pestaña \\
Ver columnas muy a la derecha & Barra de desplazamiento horizontal, o \tecla{F5} y la dirección \\
\bottomrule
\end{tabularx}

\begin{nota}
En \hoja{ENTRADA} y otras hojas, los títulos de arriba y las primeras columnas están \textbf{fijos}: no se mueven al desplazarse.
Así siempre sabe en qué contenedor y en qué columna está.
\end{nota}

\section{Escribir, corregir y borrar}

\begin{tabularx}{\linewidth}{@{}l X@{}}
\toprule
\textbf{Para\dots} & \textbf{Haga esto} \\ \midrule
Escribir en una celda vacía & Selecciónela, escriba y pulse \tecla{Enter}. \\
Reemplazar lo que hay en una celda & Selecciónela y escriba encima. Lo anterior se sustituye al pulsar \tecla{Enter}. \\
Corregir solo una parte & Selecciónela y pulse \tecla{F2} (o doble clic). Mueva el cursor con las flechas y corrija. \\
Arrepentirse mientras escribe & \tecla{Esc}: la celda queda como estaba. \\
Borrar el contenido & Seleccione la celda, o varias, y pulse \tecla{Supr}. \\
\textbf{Deshacer lo último} & \tecla{Ctrl}+\tecla{Z}. Puede pulsarlo varias veces. \\
Seleccionar varias celdas & Haga clic en la primera y arrastre el ratón; o mantenga \tecla{Mayús} y use las flechas. \\
\bottomrule
\end{tabularx}

\begin{atencion}
\textbf{Nunca elimine filas ni columnas} (clic derecho $\rightarrow$ «Eliminar»). Las fórmulas dependen de la posición de cada fila.
Para quitar datos, selecciónelos y pulse \tecla{Supr}. El libro está protegido para impedirlo, pero conviene saberlo.
\end{atencion}

\section{Números, fechas y decimales}

\begin{itemize}
  \item \textbf{Números}: escriba solo cifras, sin unidades ni símbolos: \escriba{5000}, no «5000 kg»; \escriba{2.30} o \escriba{2,30}
        (vea la nota), no «2,30 USD». El libro ya sabe en qué unidad va cada columna.
  \item \textbf{Fechas}: escriba día/mes/año: \escriba{27/10/2026}. Si al pulsar \tecla{Enter} la fecha queda alineada a la
        \textbf{derecha}, Excel la reconoció. Si queda a la izquierda, la tomó como texto: corríjala.
  \item \textbf{Porcentajes}: escriba \escriba{3\%} o \escriba{0.03}, que son lo mismo.
\end{itemize}

\begin{nota}
\textbf{Punto o coma decimal.} Depende de la configuración de Windows. En las imágenes de esta guía los números se ven
con coma de miles y punto decimal: \texttt{\$2,806.24}. En su computadora pueden verse como \texttt{2.806,24~\$}. Es el mismo número.
Use el separador decimal que use su Excel: si escribe \escriba{2,30} y Excel lo toma como texto (alineado a la izquierda),
escriba \escriba{2.30}, o al revés.
\end{nota}

\section{Listas desplegables, notas de ayuda y avisos}

\begin{itemize}
  \item \textbf{Listas desplegables}: algunas celdas (Transitaria, Tipo, Estado\dots) solo admiten ciertos valores. Al seleccionarlas aparece
        una flecha \recuadro{7}: púlsela y elija. También puede escribir el valor exacto.
  \item \textbf{Notas de ayuda}: los títulos con un triángulo rojo en la esquina \recuadro{8} tienen una explicación. Ponga el ratón encima.
  \item \textbf{Valor no válido}: si escribe algo no permitido (texto en una celda de números, una fecha imposible, un valor que no está
        en la lista), Excel muestra este aviso:
\end{itemize}

\begin{center}
\begin{tikzpicture}
\node[draw=gris,fill=white,text width=8.5cm,inner sep=8pt,font=\small,drop shadow] (v) {
  \textbf{Número no válido}\\[4pt]
  Escriba un número mayor o igual que 0, sin texto.\\[6pt]
  \hfill\fbox{\footnotesize Reintentar}\;\fbox{\footnotesize Cancelar}\;\fbox{\footnotesize Ayuda}};
\end{tikzpicture}
\end{center}
Pulse \boton{Reintentar} para corregir lo que escribió, o \boton{Cancelar} para dejar la celda como estaba.

\begin{itemize}
  \item \textbf{Hoja protegida}: si intenta escribir en una celda de fórmula (letra negra o verde), Excel avisa:
  \emph{«La celda o el gráfico que intenta cambiar está en una hoja protegida»}. Pulse \boton{Aceptar}: no pasó nada.
  Esa celda se calcula sola. Escriba solo en las celdas amarillas.
\end{itemize}

\section{Abrir y guardar}

\begin{pasos}
\paso Abra el archivo con doble clic.
\paso Si aparece una barra amarilla que dice \textbf{VISTA PROTEGIDA}, típica de archivos recibidos por correo, pulse
      \boton{Habilitar edición}. Si no lo hace, no podrá escribir.
\paso Para \textbf{guardar}, pulse \tecla{Ctrl}+\tecla{G} (Excel en español) o \tecla{Ctrl}+\tecla{S} (Excel en inglés), o el icono
      del disquete. Guarde cada vez que termine un contenedor.
\paso Para hacer una \textbf{copia con otro nombre}: \menu{Archivo $\rightarrow$ Guardar como}. Elija la carpeta y escriba el nombre,
      por ejemplo \escribal{Rentabilidad\_2026-10.xlsx}.
\end{pasos}

\begin{atencion}
Después de un «Guardar como», Excel queda trabajando con el archivo \emph{nuevo}. Si la copia era solo un respaldo, cierre y vuelva a
abrir su archivo maestro, o guarde la copia como se explica en la sección~\ref{sec:respaldo}.
\end{atencion}

\section{Hacer una copia de práctica}\label{sec:practica}

Antes de los ejercicios del capítulo~\ref{cap:taller}, haga una copia para practicar sin miedo:

\begin{pasos}
\paso Abra \libro.
\paso \menu{Archivo $\rightarrow$ Guardar como} y escriba \escribal{PRACTICA\_Rentabilidad.xlsx}.
\paso Trabaje en \texttt{PRACTICA\_Rentabilidad.xlsx}. El original queda intacto.
\end{pasos}

\section{Resumen de teclas útiles}

\begin{center}
\begin{tabular}{ll@{\hspace{2cm}}ll}
\toprule
\tecla{Enter} & confirmar y bajar & \tecla{Ctrl}+\tecla{Z} & deshacer \\
\tecla{Tab} & confirmar e ir a la derecha & \tecla{Ctrl}+\tecla{G} / \tecla{S} & guardar \\
\tecla{Esc} & cancelar lo que escribe & \tecla{Ctrl}+\tecla{C} / \tecla{V} & copiar / pegar \\
\tecla{F2} & corregir dentro de la celda & \tecla{F5} & ir a una celda \\
\tecla{Supr} & borrar el contenido & \tecla{Ctrl}+\tecla{Inicio} & ir al principio \\
\bottomrule
\end{tabular}
\end{center}

\begin{consejo}
Para copiar datos desde otro archivo, use \textbf{Pegado especial $\rightarrow$ Valores}: clic derecho $\rightarrow$ «Opciones de pegado»
$\rightarrow$ el icono «123». Así no se arrastran colores ni formatos ajenos.
\end{consejo}

% ============================================================================
\chapter{Conozca el libro, hoja por hoja}\label{cap:conozca}
% ============================================================================

Este capítulo es un recorrido. Todavía no hay que escribir nada: abra el libro y vaya mirando cada hoja a medida que la lee.

\section{La regla de oro: los colores}

\begin{center}
\begin{tabularx}{0.95\linewidth}{@{}lXc@{}}
\toprule
\textbf{Aspecto} & \textbf{Qué es} & \textbf{¿Escribo?} \\ \midrule
\entradaejemplo{} fondo amarillo, letra azul & Dato que usted escribe & \textbf{\color{verde}SÍ} \\
\textbf{1.234} letra negra & Fórmula: se calcula sola & \textbf{\color{rojo}NO} \\
\textcolor{enlace}{1.234} letra verde & Dato traído de otra hoja & \textbf{\color{rojo}NO} \\
\colorbox{azul}{\textcolor{white}{\bfseries Título}} fondo azul & Títulos de columnas & \textbf{\color{rojo}NO} \\
\textcolor{red}{(\$114)} rojo entre paréntesis & Número negativo: pérdida & \textbf{\color{rojo}NO} \\
\textbf{-} un guion & Cero & --- \\
\bottomrule
\end{tabularx}
\end{center}

\begin{nota}
Los colores se repiten en todas las hojas. Si duda si puede escribir en una celda, mire su color: \textbf{solo en las amarillas}.
\end{nota}

% ---------------------------------------------------------------- INICIO
\section{\hoja{INICIO}: la portada y el tablero de control}

Es la primera hoja que verá al abrir el libro.

\captura{inicio_indice}{Hoja INICIO: índice de hojas}
\begin{itemize}
  \item \recuadro{1} \textbf{Nombres de las hojas}, en azul subrayado: son \textbf{enlaces}. Haga clic en uno y Excel lo lleva a esa hoja.
  \item \recuadro{2} \textbf{Cuándo se usa} cada hoja.
\end{itemize}

\captura{inicio_rutina}{Hoja INICIO: su rutina de trabajo y la leyenda de colores}
\recuadro{1} Su \textbf{rutina de trabajo}: el resumen de los cuatro momentos de la sección anterior.

\captura{inicio_controles}{Hoja INICIO: controles automáticos}
\recuadro{1} \textbf{Estado del libro}. Revíselo siempre antes de dar un mes por terminado:
\begin{itemize}
  \item \textbf{Contenedores registrados}, \textbf{con alertas} y \textbf{abiertos} (pendientes de cerrar).
  \item \textbf{Cuadre de regiones} y \textbf{cuadre mensual}: deben decir \textcolor{verde}{\textbf{OK}}. Si dicen \textcolor{rojo}{\textbf{REVISAR}},
        consulte el capítulo~\ref{cap:problemas}.
  \item \textbf{Conciliación de comercialización}: debe decir \textcolor{verde}{\textbf{CUADRA}} después del cierre de mes.
  \item \textbf{Utilidad antes de impuesto acumulada}: con los datos de ejemplo, \vUAItotal.
\end{itemize}

% ---------------------------------------------------------------- ENTRADA
\section{\hoja{ENTRADA}: donde registra cada contenedor}

Es la hoja que usará casi todos los días. \textbf{Cada fila es un contenedor}. Hay espacio para 300, de la fila 6 a la 305.

\captura{entrada_izq}{Hoja ENTRADA, primera parte (columnas A a L)}
\begin{itemize}
  \item \textbf{A – Nº}: se numera solo al escribir el ID.
  \item \recuadro{1} \textbf{B – ID contenedor}: el número del contenedor o del BL. \textbf{Es obligatorio}: una fila sin ID no se calcula.
  \item \recuadro{2} \textbf{C – Fecha de arribo}, \textbf{D – Transitaria} (lista: ENCI, A.V, Otra), \textbf{E – Tipo} (lista) y
        \textbf{F – Nº de bultos} (informativo).
  \item \recuadro{3} \textbf{G, H, I – kg por región}: cuántos kilogramos van a Occidente, a Centro y a Oriente. Salen del manifiesto o
        del desglose por destinatario. \textbf{Son el dato más importante}: de ellos dependen los ingresos.
  \item \recuadro{4} \textbf{J a U – columnas «(opcional)»}: los datos reales de cada etapa (viajes, litros, dieta, jornales). Si las deja vacías,
        el libro usa la \textbf{norma técnica}, el valor estándar de \hoja{PARAMETROS}. Por ejemplo, 60 litros por viaje al puerto.
\end{itemize}

\captura{entrada_der}{Hoja ENTRADA, parte final (columnas V a AC)}
\begin{itemize}
  \item \textbf{V–W}: otros gastos directos del contenedor y su detalle. \textbf{X}: otros ingresos (entregas a domicilio, sobrepeso\dots).
  \item \recuadro{1} \textbf{Y – Fecha de fin de la distribución} y \textbf{Z – Estado}: «Abierto» mientras falten datos; «Cerrado» al completarlos.
        \textbf{AA – Observaciones}: anote aquí los números de los documentos de soporte.
  \item \recuadro{2} \textbf{AB – Alertas} y \textbf{AC – Utilidad}: \textbf{se calculan solas}. Le dicen al instante si algo está mal y cuánto deja
        el contenedor. Aquí el contenedor 3 tiene la alerta «\vAlertaTres».
\end{itemize}

La lista completa de columnas, con el documento de donde sale cada dato, está en el anexo~\ref{anx:columnas}.

% ---------------------------------------------------------------- PARAMETROS
\section{\hoja{PARAMETROS}: normas técnicas y tasas}

Datos estables que se configuran una vez y casi no cambian.

\captura[0.8\linewidth]{parametros}{Hoja PARAMETROS (columnas B a D; la E trae la explicación de cada dato)}
\begin{itemize}
  \item \recuadro{1} \textbf{Generales}: nombre de la empresa, primer mes del período (\texttt{Fecha\_Inicio}), forma de repartir los gastos
        fijos (por kg, recomendado), base de la comisión variable, tolerancia del combustible (10\,\%) y mes de corte (opcional).
  \item \recuadro{2} \textbf{Normas técnicas}: litros por viaje de cada tramo, kilómetros al puerto, capacidad de los camiones y
        rendimiento del desagrupe (kg que desagrupa un trabajador en una jornada).
  \item Más abajo: los \textbf{porcentajes sobre el ingreso} (reclamaciones y banca) y los \textbf{tributos}.
\end{itemize}
En el libro, la columna E explica cada dato y dice si es un valor ILUSTRATIVO que hay que sustituir.

% ---------------------------------------------------------------- PRECIOS
\section{\hoja{PRECIOS}: tarifas y precios con fecha}

Aquí están los precios que cambian con el tiempo. La hoja es ancha, así que se muestra en dos imágenes.

\captura{precios}{Hoja PRECIOS, columnas A a I}
\begin{itemize}
  \item \recuadro{1} \textbf{Vigente desde}: la fecha desde la que rigen los precios de esa fila.
  \item \recuadro{2} Tarifas por kg de cada región, diésel, cita, dieta y pago por jornal.
\end{itemize}

\captura{precios_b}{Hoja PRECIOS, columnas J a O}
\recuadro{1} Servicios de las transitarias, comisiones de los gestores y desgaste del camión. \recuadro{2} El documento de soporte del cambio de precio.

\begin{atencion}
\textbf{Cuando un precio cambia, no borre la fila anterior: agregue una fila nueva debajo} con la nueva fecha. Cada contenedor usa
los precios vigentes en su fecha de arribo, así que los contenedores anteriores conservan su costo real. Lo practicará en el ejercicio~4.
\end{atencion}

% ---------------------------------------------------------------- ACTIVOS
\section{\hoja{ACTIVOS}: el camión y los equipos}

\captura{activos}{Hoja ACTIVOS, columnas A a J}
\begin{itemize}
  \item \recuadro{1} Lo que usted escribe: descripción, tipo, fecha de alta, valor de compra, valor residual (lo que valdrá al final de su vida
        útil) y vida útil en años.
  \item \recuadro{2} Lo que calcula el libro: depreciación anual y mensual, y fecha en que termina de depreciarse.
\end{itemize}

\captura{activos_b}{Hoja ACTIVOS, columnas K a O}
\recuadro{1} Fecha de baja (si se vende o se retira), mantenimiento preventivo mensual y seguros y licencias anuales.

\begin{nota}
La \textbf{depreciación} es el desgaste del camión expresado en dinero. Si costó 45\,000 USD, al final valdrá 5\,000 y dura 8 años,
pierde $(45\,000-5\,000)/8 = 5\,000$ USD al año, es decir, 416,67 USD al mes. Es un gasto real aunque no salga dinero de la caja:
cuando haya que reponer el camión, habrá que tener ese dinero.
\end{nota}

% ---------------------------------------------------------------- COSTOS FIJOS
\section{\hoja{COSTOS\_FIJOS}: los gastos de cada mes}

Los \textbf{gastos fijos} son los que se pagan aunque no llegue ningún contenedor: salarios de oficina, alquiler, electricidad,
servicios contratados de chofer y ayudante, depreciación\dots

\captura{costos_fijos}{Hoja COSTOS\_FIJOS: conceptos, valor base y meses}
\begin{itemize}
  \item \textbf{C – Centro de costo}: GENERAL, si es de toda la empresa, u OCC/CEN/ORI, si es de una región.
        Los gastos de una región solo se cargan a los kg de esa región.
  \item \recuadro{1} \textbf{¿Salario de trabajador propio?}: «Sí» para los salarios de los trabajadores de la empresa. Con eso el libro calcula
        la seguridad social (14\,\%) y el impuesto por la fuerza de trabajo (5\,\%). «No» para servicios contratados a terceros.
  \item \recuadro{2} \textbf{Valor base mensual}: lo que se paga normalmente cada mes.
  \item \recuadro{3} \textbf{Una columna por mes} (24 meses). Cada mes copia el valor base. Si un mes el gasto real fue distinto, escriba el
        importe real en la celda de ese mes.
\end{itemize}

\captura{costos_fijos_tot}{Hoja COSTOS\_FIJOS: tributos sobre salarios y totales}
\recuadro{1} Los tributos sobre salarios se calculan solos. \recuadro{2} Los totales por centro de costo y el
\textbf{total de costos fijos del mes}: \vFijosMes{} en el ejemplo.

% ---------------------------------------------------------------- COMERCIAL
\section{\hoja{COMERCIAL}: los 21 gestores}

\captura{comercial}{Hoja COMERCIAL: datos del mes y lista de gestores}
\begin{itemize}
  \item \recuadro{1} \textbf{Mes a liquidar}: el mes del que va a calcular el pago.
  \item Lista de los 21 gestores: 5 del grupo central y 16 provinciales, uno por territorio.
  \item \recuadro{2} \textbf{F – Pago fijo mensual} (igual para todos; viene de la celda \celda{C8}) y \textbf{G – Base variable del mes}: lo que
        gestionó cada uno, en USD cobrados o en kg, según la base elegida en \hoja{PARAMETROS}.
  \item \textbf{H} e \textbf{I}: pago variable y total a pagar, calculados.
\end{itemize}

\captura{comercial_conc}{Hoja COMERCIAL: conciliación}
\recuadro{1} \textbf{Conciliación}: compara lo que declaran los gestores con lo que dicen los contenedores registrados. Si coinciden,
dice \textcolor{verde}{\textbf{CUADRA}}. Si no, \textcolor{rojo}{\textbf{REVISAR}}. Lo practicará en el ejercicio~5.

% ---------------------------------------------------------------- CALCULO
\section{\hoja{CALCULO}: las cuentas, a la vista (solo consulta)}

Aquí el libro hace todas las cuentas de cada contenedor, en unas 95 columnas. \textbf{No se escribe nada}, pero puede consultarla para
entender de dónde sale un resultado. Las imágenes muestran solo algunas columnas.

\captura{calculo}{Hoja CALCULO: algunas columnas clave}
\recuadro{1} El contenedor 3: sus ingresos, el costo de cada etapa, los costos fijos que le tocan y la utilidad.

\captura{calculo_region}{Hoja CALCULO: utilidad por región, control y alertas}
\recuadro{1} El contenedor 3 gana en Occidente y en Centro, pero pierde en Oriente. El \textbf{control de cuadre} vale 0 (aparece
como «-»): la suma de las tres regiones es igual a la utilidad del contenedor.

% ---------------------------------------------------------------- RESUMEN
\section{\hoja{RESUMEN\_MES}: el resultado de cada mes}

\captura{resumen}{Hoja RESUMEN\_MES: columnas principales (se ocultaron algunas para que se lea mejor)}
\recuadro{1} Una fila por mes. \recuadro{2} \textbf{Utilidad antes de impuesto} del mes: \vUAIsep{} en septiembre y \vUAIoct{} en octubre,
con los datos de ejemplo. Noviembre no carga gastos porque todavía no ha ocurrido: es posterior al «mes de corte».

\captura{resumen_pe}{Hoja RESUMEN\_MES: indicadores por kg y punto de equilibrio}
\recuadro{1} \textbf{Punto de equilibrio} y \textbf{margen de seguridad}. En septiembre bastaban \vPEsep{} kg (\vPEcontSep{} contenedores)
para cubrir todos los gastos. El capítulo~\ref{cap:resultados} explica cada indicador.

% ---------------------------------------------------------------- DASHBOARD
\section{\hoja{DASHBOARD}: el panel para informar}

\captura{dashboard_kpi}{DASHBOARD: período e indicadores}
\recuadro{1} Elija el \textbf{período}: mes «desde» y mes «hasta», con la flecha de la lista. \recuadro{2} Utilidad antes de impuesto,
impuesto estimado y utilidad neta del período. Los demás indicadores se explican en el capítulo~\ref{cap:resultados}.

\captura{dashboard_tablas}{DASHBOARD: estructura de costos y rentabilidad por región}
\recuadro{1} \textbf{Rentabilidad por región}. Con los datos de ejemplo, Occidente gana \vUAIocc, Centro \vUAIcen{} y Oriente
\textcolor{red}{\vUAIori}, es decir, pierde dinero.

\captura{dashboard_graficos}{DASHBOARD: gráficos}

% ---------------------------------------------------------------- SIMULADOR
\section{\hoja{SIMULADOR}: «¿qué pasaría si\dots?»}

Sirve para probar escenarios \textbf{sin tocar los datos reales}: subir una tarifa, un aumento del diésel, recibir más contenedores\dots

\captura{simulador_sup}{SIMULADOR: supuestos}
\recuadro{1} \textbf{Valor actual}: viene de los datos reales del libro. \recuadro{2} \textbf{Valor a simular}: escriba aquí para probar.
\recuadro{3} \textbf{Valor usado}: el simulado, si escribió uno; si no, el actual.

\captura{simulador_res}{SIMULADOR: resultado de un contenedor típico y del mes}
\recuadro{1} Utilidad de un contenedor típico por región. \recuadro{2} \textbf{Tarifa mínima de equilibrio} de cada región: por debajo de ese
precio por kg, la región pierde dinero. Oriente necesita al menos \vSimTminOri{} por kg; Occidente, solo \vSimTminOcc.
\recuadro{3} Punto de equilibrio: contenedores por mes necesarios para no perder (\vSimPE).

\captura{simulador_sens}{SIMULADOR: sensibilidad a la tarifa y al diésel}
Cada casilla es la utilidad del mes con una variación de las tarifas (filas) y un precio del diésel (columnas). \recuadro{1} es la situación
actual: tarifas sin cambio y diésel a 2,00 USD/L, \vSimGrid. Si el diésel sube a 2,50 USD/L, la utilidad baja a \vSimGridMas. Las casillas
rosadas son pérdidas.

\captura[0.6\linewidth]{simulador_vol}{SIMULADOR: utilidad según la cantidad de contenedores al mes}

% ---------------------------------------------------------------- otras
\section{Las demás hojas}
\begin{itemize}
  \item \hoja{GUIA\_USO}: un resumen de esta guía dentro del propio libro.
  \item \hoja{ANALISIS}: el diagnóstico del negocio, la metodología y las propuestas de mejora.
  \item \hoja{LISTAS}: las opciones de las listas desplegables. No se toca.
\end{itemize}

% ============================================================================
\chapter{Taller práctico: 8 ejercicios guiados}\label{cap:taller}
% ============================================================================

En este taller hará, una vez y de principio a fin, todo lo que hará después en su trabajo real. Cada ejercicio indica
\textbf{qué escribir y en qué celda}, y al final trae un cuadro verde con \textbf{los números exactos que debe ver}.
Si coinciden, lo hizo bien.

\begin{atencion}
\textbf{Haga los ejercicios sobre una copia de práctica} (sección~\ref{sec:practica}): \texttt{PRACTICA\_Rentabilidad.xlsx}.
Hágalos \textbf{en orden}: cada ejercicio parte de cómo quedó el anterior. Si se equivoca, use \tecla{Ctrl}+\tecla{Z}.
Si se pierde del todo, vuelva a copiar el archivo original y empiece de nuevo.
\end{atencion}

\begin{center}
\begin{tabularx}{\linewidth}{@{}clX@{}}
\toprule
\textbf{Nº} & \textbf{Ejercicio} & \textbf{Qué aprende} \\ \midrule
1 & Comprobar el estado del libro & Leer los controles de \hoja{INICIO} \\
2 & Registrar un contenedor que llega & El primer momento de cada contenedor \\
3 & Completar los datos reales y cerrarlo & El segundo momento; cómo nace una alerta \\
4 & Sube el precio del diésel & Agregar una vigencia de precios sin alterar el pasado \\
5 & Cierre del mes de octubre & Gastos reales, liquidación de los gestores y conciliación \\
6 & Informe del período & Usar el \hoja{DASHBOARD} \\
7 & ¿Y si subimos la tarifa de Oriente? & Usar el \hoja{SIMULADOR} \\
8 & Dejar el libro limpio para empezar & Borrar los ejemplos de forma segura \\
\bottomrule
\end{tabularx}
\end{center}

% ============================================================================
\section{Ejercicio 1. Comprobar el estado del libro}

\begin{ficha}{Ejercicio 1 · Duración: 3 minutos}
\textbf{Objetivo:} aprender a leer el tablero de control antes de trabajar.\\
\textbf{Hoja:} \hoja{INICIO}
\end{ficha}

\begin{pasos}
\paso Abra \texttt{PRACTICA\_Rentabilidad.xlsx}. Si aparece la barra amarilla de \textbf{Vista protegida}, pulse \boton{Habilitar edición}.
\paso Haga clic en la pestaña \hoja{INICIO}.
\paso Baje hasta el bloque \textbf{Estado del libro (controles automáticos)}, alrededor de la fila 35.
\paso Lea cada línea y compárela con el cuadro verde.
\end{pasos}

\captura{inicio_controles}{Ejercicio 1: el estado inicial del libro}

\begin{resultado}
Contenedores registrados \textbf{4}, con alertas \textbf{\vAlertasIni}, abiertos \textbf{\vAbiertosIni};
cuadres \textbf{OK} y \textbf{OK}; conciliación \textbf{CUADRA}; utilidad acumulada \textbf{\vUAItotal}.
\end{resultado}

\begin{nota}
Haga clic en algún nombre de hoja del índice (filas 10 a 22): Excel lo lleva a esa hoja. Para volver, haga clic en la pestaña \hoja{INICIO}.
\end{nota}

% ============================================================================
\section{Ejercicio 2. Registrar un contenedor que llega}\label{sec:ej2}

\begin{ficha}{Ejercicio 2 · Duración: 5 minutos}
\textbf{Situación:} el 27 de octubre de 2026 llega el contenedor \textbf{CONT-PRUEBA-05}. Pasa por ENCI, es un 40' HC con 1\,000 bultos.
Según el manifiesto, lleva \textbf{5\,000 kg} para Occidente, \textbf{2\,500 kg} para Centro y \textbf{2\,500 kg} para Oriente.\\
\textbf{Hoja:} \hoja{ENTRADA} · \textbf{Fila:} 10, la primera vacía
\end{ficha}

\begin{pasos}
\paso Vaya a la hoja \hoja{ENTRADA}. Las filas 6 a 9 tienen los 4 contenedores de ejemplo. La primera fila libre es la \textbf{10}.
\paso Haga clic en \celda{B10} y escriba \escriba{CONT-PRUEBA-05}. Pulse \tecla{Tab}: pasa a \celda{C10}. En \celda{A10} aparece el número 5.
\paso En \celda{C10} escriba la fecha \escriba{27/10/2026} y pulse \tecla{Tab}.
\paso En \celda{D10} pulse la flecha de la lista y elija \textbf{ENCI}. Pulse \tecla{Tab}.
\paso En \celda{E10} elija \textbf{40' HC} de la lista. Pulse \tecla{Tab}.
\paso En \celda{F10} escriba \escriba{1000} y pulse \tecla{Tab}.
\paso En \celda{G10} escriba \escriba{5000}; en \celda{H10}, \escriba{2500}; en \celda{I10}, \escriba{2500}. Pulse \tecla{Tab} entre una y otra.
\paso Deje vacías las columnas J a X: todavía no tiene los datos reales.
\paso Vaya a \celda{Z10} (pulse \tecla{F5}, escriba \escriba{Z10} y \tecla{Enter}) y elija \textbf{Abierto} en la lista.
\paso Guarde con \tecla{Ctrl}+\tecla{G}.
\end{pasos}

\captura{ej2_entrada}{Ejercicio 2: el contenedor nuevo en la fila 10 (columnas A a I)}
\captura{ej2_estado}{Ejercicio 2: estado y resultado del contenedor (columnas Y a AC)}

\begin{resultado}
\begin{itemize}
  \item \recuadro{1} \celda{Z10} dice \textbf{Abierto}.
  \item \recuadro{2} \celda{AB10} (alertas) está \textbf{vacía} y \celda{AC10} (utilidad) muestra \textbf{\vEjDosUAI}.
\end{itemize}
\end{resultado}

\begin{nota}
\textbf{¿Por qué cambió la utilidad de otro contenedor?} Antes del ejercicio, el contenedor 3 (fila 8) tenía una utilidad de \vUAIcTres.
Ahora tiene \vEjDosUAItres. No es un error: los \textbf{gastos fijos de octubre} se reparten entre todos los contenedores de octubre según
sus kilogramos. Al llegar uno más, a cada uno le toca una parte menor. Por eso \textbf{el resultado de un mes es definitivo solo cuando el
mes se cierra}.
\end{nota}

% ============================================================================
\section{Ejercicio 3. Completar los datos reales y cerrar el contenedor}\label{sec:ej3}

\begin{ficha}{Ejercicio 3 · Duración: 5 minutos}
\textbf{Situación:} terminó la distribución del CONT-PRUEBA-05 el 6 de noviembre. Los vales dicen que el viaje al puerto consumió
\textbf{64 litros}, que el desagrupe en ENCI llevó \textbf{11 jornales} y que el traslado a Oriente consumió \textbf{520 litros}.
Los demás datos fueron los normales.\\
\textbf{Hoja:} \hoja{ENTRADA} · \textbf{Fila:} 10
\end{ficha}

\begin{pasos}
\paso En \celda{K10} (Litros reales al puerto) escriba \escriba{64}.
\paso En \celda{M10} (Jornales del desagrupe) escriba \escriba{11}.
\paso En \celda{S10} (Litros de traslado a Oriente) escriba \escriba{520}.
\paso Las demás columnas opcionales quedan vacías: el libro usará la norma.
\paso En \celda{Y10} escriba la fecha de fin \escriba{06/11/2026}.
\paso En \celda{Z10} cambie el estado a \textbf{Cerrado}.
\paso En \celda{AA10} anote los documentos: \escribal{Vales 1520-1528; nómina de desagrupe 45}.
\end{pasos}

\captura{ej3_reales}{Ejercicio 3: datos reales en las columnas opcionales}
\captura{ej3_estado}{Ejercicio 3: contenedor cerrado, con una alerta}

\begin{resultado}
\begin{itemize}
  \item \recuadro{2} \celda{AB10} muestra la alerta: «\textbf{\vEjTresAlerta}».
  \item \celda{AC10}: la utilidad baja a \textbf{\vEjTresUAI}.
  \item En \hoja{INICIO}, «Contenedores con alertas» pasa a \textbf{\vEjTresAlertas}.
\end{itemize}
\end{resultado}

\captura{ej3_inicio}{Ejercicio 3: el control de INICIO cuenta la nueva alerta}

\begin{nota}
\textbf{¿Qué significa la alerta?} La norma para el traslado a Oriente es de 450 litros por viaje. Se gastaron 520, un 16\,\% más, y la
tolerancia es del 10\,\%. El libro no impide registrarlo: \textbf{avisa} para que usted averigüe la causa (un desvío, una avería, una ruta
distinta, un vale mal anotado\dots).
\begin{itemize}
  \item Si el dato es correcto y está justificado, explíquelo en Observaciones.
  \item Si era un error de anotación, corríjalo. La alerta desaparece sola.
  \item Si el exceso se repite con todos los viajes, la norma está desactualizada: corríjala en \hoja{PARAMETROS} (capítulo~\ref{cap:cambios}).
\end{itemize}
\end{nota}

% ============================================================================
\section{Ejercicio 4. Sube el precio del diésel}\label{sec:ej4}

\begin{ficha}{Ejercicio 4 · Duración: 8 minutos}
\textbf{Situación:} desde el \textbf{1 de noviembre de 2026} el diésel pasa de 2,00 a \textbf{2,30 USD/L}. El 10 de noviembre llega el
contenedor \textbf{CONT-PRUEBA-06} por A.V: un 40' con 820 bultos y 4\,000 / 2\,000 / 2\,000 kg.\\
\textbf{Hojas:} \hoja{PRECIOS} y \hoja{ENTRADA}
\end{ficha}

\begin{atencion}
\textbf{No cambie el 2,00 de la fila 6 por 2,30.} Si lo hiciera, el libro recalcularía todos los contenedores anteriores con el precio nuevo
y el historial quedaría falseado. Lo correcto es \textbf{agregar una fila} con la nueva fecha.
\end{atencion}

\begin{pasos}
\paso Vaya a \hoja{PRECIOS}. La fila 6 tiene los precios vigentes desde el 01/01/2026.
\paso Seleccione \celda{B6:N6}: haga clic en \celda{B6}, mantenga \tecla{Mayús} y haga clic en \celda{N6}. Cópielo con \tecla{Ctrl}+\tecla{C}.
\paso Haga clic en \celda{B7} y pegue con \tecla{Ctrl}+\tecla{V}. La fila 7 queda igual que la 6.
\paso En \celda{A7} escriba \escriba{01/11/2026}: la fecha desde la que rige el cambio.
\paso En \celda{E7} (Precio del diésel) escriba \escriba{2.30}.
\paso En \celda{O7} anote el soporte: \escribal{Aumento del diésel (factura CUPET 0457)}.
\paso Pulse \tecla{Esc} para quitar la marca de copiado.
\paso Vaya a \hoja{ENTRADA} y registre el nuevo contenedor en la fila \textbf{11}, como en el ejercicio 2:
      \celda{B11} \escriba{CONT-PRUEBA-06}, \celda{C11} \escriba{10/11/2026}, \celda{D11} \textbf{A.V}, \celda{E11} \textbf{40'},
      \celda{F11} \escriba{820}, \celda{G11} \escriba{4000}, \celda{H11} \escriba{2000}, \celda{I11} \escriba{2000}, \celda{Z11} \textbf{Abierto}.
\end{pasos}

\captura{ej4_precios}{Ejercicio 4: la nueva vigencia en la fila 7 (se ocultaron las columnas G a N)}

Para comprobar que cada contenedor tomó su precio, mire la hoja \hoja{CALCULO}. La imagen muestra solo algunas columnas:

\captura{ej4_calculo}{Ejercicio 4: cada contenedor usa los precios de su fecha}

\begin{resultado}
\begin{itemize}
  \item Columna «Precios vigentes desde»: el CONT-PRUEBA-06 (fila 11) usa la vigencia del \textbf{01/11/2026}. Todos los demás siguen con la del 01/01/2026.
  \item Combustible al puerto del CONT-PRUEBA-06: \textbf{\vEjCuatroCombSeis} (60 L $\times$ 2,30). El CONT-PRUEBA-05 sigue con \textbf{\vEjCuatroCombCinco}
        (64 L $\times$ 2,00).
  \item La utilidad del CONT-PRUEBA-05 \textbf{no cambió}: sigue en \textbf{\vEjCuatroUAIcinco}. El historial está intacto.
  \item El CONT-PRUEBA-06 muestra una pérdida de \textcolor{red}{\textbf{\vEjCuatroUAIseis}}.
\end{itemize}
\end{resultado}

\begin{nota}
\textbf{¿Por qué pierde el CONT-PRUEBA-06?} Es, por ahora, el \textbf{único} contenedor de noviembre, así que carga solo con todos los gastos
fijos del mes. Cuando se registren los demás contenedores de noviembre, el gasto se repartirá y su resultado mejorará.
Un mes con pocos contenedores es caro: lo verá en el indicador «punto de equilibrio» (capítulo~\ref{cap:resultados}).
\end{nota}

% ============================================================================
\section{Ejercicio 5. Cierre del mes de octubre}\label{sec:ej5}

\begin{ficha}{Ejercicio 5 · Duración: 15 minutos}
\textbf{Situación:} es 1 de noviembre y toca cerrar octubre. La factura de electricidad de octubre fue de \textbf{55 USD}, no de 40.
Los gestores informan lo que cobraron en octubre. Al cobro habitual se suma el del CONT-PRUEBA-05:
\textbf{4\,000 USD} en La Habana, \textbf{2\,000} en Villa Clara y \textbf{2\,000} en Santiago de Cuba.\\
\textbf{Hojas:} \hoja{COSTOS\_FIJOS}, \hoja{COMERCIAL}, \hoja{RESUMEN\_MES}
\end{ficha}

\subsection*{Parte A: registrar un gasto fijo real distinto del habitual}
\begin{pasos}
\paso Vaya a \hoja{COSTOS\_FIJOS}. Busque la fila \textbf{10, «Electricidad y agua»}.
\paso Las columnas de los meses empiezan en G (sep-2026). Octubre es la columna \textbf{H}.
\paso Haga clic en \celda{H10}. Verá que contiene \texttt{=\$F10}: copia el valor base. Escriba encima \escriba{55} y pulse \tecla{Enter}.
\end{pasos}

\captura{ej5_costos}{Ejercicio 5A: gasto real de electricidad de octubre}

\begin{resultado}
En \hoja{RESUMEN\_MES}, «Costos fijos del mes» de octubre pasa de \vEjCincoFijosOctAntes{} a \textbf{\vEjCincoFijosOct}. Solo cambia octubre: los demás meses siguen con 40.
\end{resultado}

\subsection*{Parte B: liquidar y conciliar a los gestores}
\begin{pasos}
\paso Vaya a \hoja{COMERCIAL}. Compruebe que en \celda{C4} (Mes a liquidar) dice \textbf{oct-2026}. Si no, escriba \escriba{01/10/2026}.
\paso Baje a la \textbf{conciliación} (filas 35 a 40). Dice \textcolor{rojo}{\textbf{REVISAR}}, con una diferencia de
      \textcolor{red}{\vEjCincoDif}: los gestores declararon \vEjCincoDeclarada, pero los contenedores registrados suman \vEjCincoReal.
      Falta declarar el cobro del CONT-PRUEBA-05.
\end{pasos}

\captura{ej5_comercial_revisar}{Ejercicio 5B: la conciliación detecta una diferencia}

\begin{pasos}
\setcounter{paso}{2}
\paso Columna G, base variable del mes. Sume el cobro nuevo al que ya tenía cada gestor:
  \begin{itemize}
    \item La Habana, fila 19: \celda{G19} tenía \vEjCincoHabAntes. Escriba \escriba{\vEjCincoHab}.
    \item Villa Clara, fila 24: \celda{G24} tenía \vEjCincoVcAntes. Escriba \escriba{\vEjCincoVc}.
    \item Santiago de Cuba, fila 31: \celda{G31} tenía \vEjCincoScAntes. Escriba \escriba{\vEjCincoSc}.
  \end{itemize}
\end{pasos}

\begin{consejo}
Para sumar sin calculadora, escriba la operación en la celda: \escriba{=\vEjCincoHabAntes+4000}. Si su Excel usa coma decimal, escriba la
cifra con coma.
\end{consejo}

\captura{ej5_comercial_ok}{Ejercicio 5B: bases corregidas y conciliación cuadrada}

\begin{resultado}
\recuadro{4} La conciliación dice \textcolor{verde}{\textbf{CUADRA}}. El total a pagar a los 21 gestores pasa de \vEjCincoPagoTotalAntes{} a
\textbf{\vEjCincoPagoTotal}. Esa es la liquidación de octubre.
\end{resultado}

\subsection*{Parte C: revisar el resultado y guardar el respaldo}
\begin{pasos}
\setcounter{paso}{3}
\paso Vaya a \hoja{RESUMEN\_MES} y mire la fila de octubre.
\paso Vaya a \hoja{INICIO} y revise el estado del libro.
\paso Guarde el archivo y haga la copia de respaldo del mes (sección~\ref{sec:respaldo}).
\end{pasos}

\captura{ej5_resumen}{Ejercicio 5C: resultado de octubre después del cierre}

\begin{resultado}
\recuadro{1} Costos fijos de octubre: \textbf{\vEjCincoFijosOct}. \recuadro{2} Utilidad antes de impuesto de octubre: \textbf{\vEjCincoUAIoct}.
Noviembre, con un solo contenedor, va en \textcolor{red}{\vEjCincoUAInov}.
\end{resultado}

% ============================================================================
\section{Ejercicio 6. Informe del período}

\begin{ficha}{Ejercicio 6 · Duración: 3 minutos}
\textbf{Situación:} la dirección pide el resultado de octubre y noviembre.\\
\textbf{Hoja:} \hoja{DASHBOARD}
\end{ficha}

\begin{pasos}
\paso Vaya a \hoja{DASHBOARD}.
\paso Haga clic en \celda{C4} (Período desde), pulse la flecha y elija \textbf{oct-2026}. También puede escribir \escriba{01/10/2026}.
\paso En \celda{C5} (Período hasta) elija \textbf{nov-2026}.
\end{pasos}

\captura{ej6_dashboard}{Ejercicio 6: indicadores de octubre y noviembre}

\begin{resultado}
Contenedores: \textbf{\vEjSeisCont}. Utilidad antes de impuesto: \textbf{\vEjSeisUAI}. Margen neto: \textbf{\vEjSeisMargen}.
Los gráficos de más abajo se actualizan solos.
\end{resultado}

\begin{consejo}
Para presentar el informe: \menu{Archivo $\rightarrow$ Imprimir} con la hoja \hoja{DASHBOARD} activa, o \menu{Archivo $\rightarrow$
Exportar $\rightarrow$ PDF}. La hoja ya está preparada para salir en horizontal y ajustada al ancho de la página.
\end{consejo}

% ============================================================================
\section{Ejercicio 7. ¿Y si subimos la tarifa de Oriente?}

\begin{ficha}{Ejercicio 7 · Duración: 5 minutos}
\textbf{Situación:} Oriente casi no deja ganancia. La dirección pregunta qué pasaría si su tarifa subiera de 0,80 a \textbf{0,95 USD/kg}.\\
\textbf{Hoja:} \hoja{SIMULADOR}
\end{ficha}

\begin{pasos}
\paso Vaya a \hoja{SIMULADOR}. Mire la fila 11, «Tarifa Oriente».
\paso En \celda{D11} (Valor a simular) escriba \escriba{0.95}.
\paso Mire el resultado más abajo, en las filas 38 a 46.
\paso Cuando termine, \textbf{borre \celda{D11}} con \tecla{Supr} para volver al valor real.
\end{pasos}

\captura{ej7_simulador}{Ejercicio 7: tarifa de Oriente simulada}
\captura{ej7_resultado}{Ejercicio 7: nuevo resultado}

\begin{resultado}
\begin{itemize}
  \item Utilidad de Oriente por contenedor: de \vSimUAIori{} a \textbf{\vEjSieteUAIori}. Margen de Oriente: de \vSimMgOri{} a \textbf{\vEjSieteMgOri}.
  \item Utilidad del mes: de \vSimUAImes{} a \textbf{\vEjSieteUAImes}.
\end{itemize}
\end{resultado}

\begin{nota}
El simulador \textbf{nunca} modifica los datos reales. Puede probar tantos escenarios como quiera. Recuerde borrar los valores simulados al terminar
para no confundirse la próxima vez.
\end{nota}

% ============================================================================
\section{Ejercicio 8. Dejar el libro limpio para empezar}\label{sec:ej8}

\begin{ficha}{Ejercicio 8 · Duración: 5 minutos}
\textbf{Objetivo:} aprender a borrar datos sin dañar el libro. Es lo que hará en el archivo real antes de cargar sus datos.\\
\textbf{Hojas:} \hoja{ENTRADA}, \hoja{COMERCIAL}
\end{ficha}

\begin{pasos}
\paso Para el ejercicio, cierre la copia de práctica sin guardar y vuelva a hacer una copia del original: así parte de los 4 contenedores de ejemplo.
\paso En \hoja{ENTRADA}, haga clic en \celda{B6} y arrastre hasta \celda{AA9}, de modo que queden seleccionadas las 4 filas de ejemplo.
\paso Pulse \tecla{Supr}. \textbf{No} use «Eliminar filas».
\paso En \hoja{COMERCIAL}, seleccione \celda{G12:G32} (base variable de los gestores) y pulse \tecla{Supr}.
\paso Vaya a \hoja{INICIO} y mire el estado del libro.
\end{pasos}

\captura{ej8_vacio}{Ejercicio 8: ENTRADA sin los ejemplos}
\captura{ej8_inicio}{Ejercicio 8: el estado del libro vacío}

\begin{resultado}
Contenedores registrados \textbf{0}, alertas \textbf{0}, cuadres \textbf{OK}, conciliación \textbf{\vEjOchoConc} y utilidad acumulada
\textcolor{red}{\textbf{\vEjOchoUAI}}.
\end{resultado}

\begin{nota}
\textbf{¿Por qué la utilidad es negativa si no hay contenedores?} El período empieza en septiembre de 2026 (\texttt{Fecha\_Inicio}), y ese mes ya
tiene gastos fijos (\vEjOchoFijos) aunque no se haya registrado ningún contenedor. Es lo correcto: los gastos fijos se pagan haya o no carga.
Al cargar sus datos reales pondrá su propia fecha de inicio (capítulo~\ref{cap:preparar}).
\end{nota}

\begin{consejo}
\textbf{¡Felicidades!} Si llegó hasta aquí con todos los cuadros verdes coincidiendo, ya sabe hacer todo el trabajo del libro.
El siguiente paso es prepararlo con los datos reales de la empresa.
\end{consejo}

% ============================================================================
\chapter{Preparar el libro con sus datos reales}\label{cap:preparar}
% ============================================================================

Este capítulo se hace \textbf{una sola vez}, antes de empezar a registrar contenedores reales. Reserve entre 2 y 4 horas y tenga
a mano los documentos de la lista siguiente.

\section{Reúna los documentos}

\begin{center}
\begin{tabularx}{\linewidth}{@{}>{\raggedright\arraybackslash}p{4.6cm}>{\raggedright\arraybackslash}X l@{}}
\toprule
\textbf{Dato} & \textbf{Dónde se obtiene} & \textbf{Hoja} \\ \midrule
Tarifas por kg que cobra la empresa & Contratos con las agencias emisoras o lista de precios & \hoja{PRECIOS} \\
Precio del diésel & Última factura o vale de CUPET & \hoja{PRECIOS} \\
Cita del puerto y dieta del chofer & Comprobantes de pago & \hoja{PRECIOS} \\
Pago por jornal (transitaria y centros) & Nómina o contrato de los trabajadores del desagrupe & \hoja{PRECIOS} \\
Cargos de ENCI y A.V por contenedor & Facturas o tarifas escritas de cada transitaria & \hoja{PRECIOS} \\
Comisión variable de los gestores & Contratos de los gestores & \hoja{PRECIOS} \\
Valor del camión y de los equipos & Registro de activos fijos (contabilidad) & \hoja{ACTIVOS} \\
Mantenimiento, seguros, licencias & Facturas del último año & \hoja{ACTIVOS} \\
Salarios, alquiler, electricidad, comunicaciones\dots & Nóminas y facturas de un mes típico & \hoja{COSTOS\_FIJOS} \\
Pago a choferes y ayudantes contratados & Contratos de servicio por región & \hoja{COSTOS\_FIJOS} \\
Nombres y territorio de los 21 gestores & Plantilla de comercialización & \hoja{COMERCIAL} \\
Tasas de impuestos & Contador / ONAT & \hoja{PARAMETROS} \\
\bottomrule
\end{tabularx}
\end{center}

\section{Cree su archivo maestro}

\begin{pasos}
\paso Copie el archivo original \libro{} y guárdelo como \escribal{Rentabilidad\_\allowbreak Paqueteria\_\allowbreak MAESTRO.xlsx} en una carpeta fija, por ejemplo
      \texttt{Documentos\textbackslash Rentabilidad}.
\paso Guarde el original sin tocar en otra carpeta. Si algún día algo sale muy mal, podrá empezar de nuevo desde él.
\paso A partir de ahora trabaje \textbf{siempre} en el archivo maestro. No trabaje en dos copias a la vez.
\end{pasos}

\section{Hoja PARAMETROS}

\captura[0.8\linewidth]{parametros}{PARAMETROS: los valores que debe revisar}

\begin{pasos}
\paso \celda{C6} \textbf{Empresa}: escriba el nombre de la empresa.
\paso \celda{C7} \textbf{Fecha\_Inicio}: el \textbf{primer día del primer mes} que va a registrar, por ejemplo \escriba{01/11/2026}.
      Si quiere cargar también meses anteriores con datos que ya tiene, ponga el primero de ellos.
\paso \celda{C8} \textbf{Metodo\_Prorrateo}: deje \textbf{KG}. Los gastos fijos se reparten según el peso, que es lo más justo.
\paso \celda{C9} \textbf{Base\_Comision}: \textbf{\% DEL INGRESO} si a los gestores se les paga un porcentaje de lo cobrado, o
      \textbf{USD POR KG} si se les paga un importe por kilogramo.
\paso \celda{C10} \textbf{Tolerancia\_Comb}: 10\,\% es razonable. Con más tolerancia habrá menos alertas, y con menos, más control.
\paso \celda{C11} \textbf{Mes\_Corte}: déjelo \textbf{vacío}.
\paso \textbf{Normas técnicas} (filas 13 a 25): sustituya los valores ilustrativos por los medidos (sección~\ref{sec:medir}).
\paso \textbf{Reclamaciones y banca} (filas 27 y 28): el porcentaje histórico de pérdidas y reclamaciones y lo que cobra el banco. Si no aplica, \escriba{0}.
\paso \textbf{Tributos} (filas 30 a 34): confírmelos con el contador. Los valores del libro son los del régimen general de las MIPYMES en 2026.
\end{pasos}

\subsection{Cómo medir las normas técnicas}\label{sec:medir}

Las normas son el consumo «normal» de cada tramo. Mídalas con 2 o 3 contenedores reales:

\begin{center}
\begin{tabularx}{\linewidth}{@{}l>{\raggedright\arraybackslash}X@{}}
\toprule
\textbf{Norma} & \textbf{Cómo se mide} \\ \midrule
Litros por viaje al puerto & Litros de los vales de cada viaje ÷ número de viajes. Promedie. \\
Km por viaje al puerto & Odómetro al volver $-$ odómetro al salir, en la hoja de ruta. \\
Rendimiento del desagrupe & kg del contenedor ÷ jornales usados. Ejemplo: 10\,000 kg con 10 jornales = 1\,000 kg/jornal. \\
Capacidad del camión & kg que el camión contratado carga realmente por viaje. \\
Litros de traslado y de distribución & Litros de los vales de cada tramo y región ÷ viajes. \\
\bottomrule
\end{tabularx}
\end{center}

\section{Hoja PRECIOS}

\begin{pasos}
\paso En la fila 6, cambie \celda{A6} por una fecha \textbf{igual o anterior} a su Fecha\_Inicio, por ejemplo \escriba{01/01/2026}.
\paso Escriba los precios reales de cada columna, de la B a la N. La columna O es para anotar el documento de soporte.
\paso Si hay servicios de las transitarias por contenedor, escríbalos en J (ENCI) y K (A.V). Si no cobran aparte, \escriba{0}.
\paso Escriba la comisión en la columna que corresponda a su Base\_Comision: L, si es un porcentaje; M, si es en USD por kg.
\end{pasos}

\section{Hoja ACTIVOS}

\begin{pasos}
\paso Fila 6: cambie los datos del camión de ejemplo por los de su camión real: chapa, fecha de alta, valor de adquisición, valor residual y
      vida útil. Pida estos datos al contador.
\paso Fila 7: escriba sus equipos reales, o borre la fila con \tecla{Supr} si no los quiere incluir.
\paso Columnas L y M: mantenimiento preventivo mensual y seguros, licencias e inspecciones anuales del camión.
\paso Si tiene más activos, use las filas 8 a 15.
\end{pasos}

\section{Hoja COSTOS\_FIJOS}

\captura[0.85\linewidth]{costos_fijos}{COSTOS\_FIJOS: columnas E (¿salario propio?) y F (valor base)}

\begin{pasos}
\paso Recorra los conceptos de las filas 6 a 27 y escriba en la columna \textbf{F} lo que se paga en un mes normal.
\paso En la columna \textbf{E}, «Sí» para los salarios de trabajadores propios y «No» para todo lo demás (alquiler, servicios contratados, etc.).
\paso La fila 8 (comercialización fija) y las filas 15 a 17 (depreciación, mantenimiento y seguros) se llenan solas desde \hoja{COMERCIAL} y \hoja{ACTIVOS}.
\paso Si un concepto no existe, ponga \escriba{0}. Si falta alguno, use una fila «Otros gastos fijos» y escriba su nombre en la columna B
      (los nombres de los conceptos son editables).
\end{pasos}

\section{Hoja COMERCIAL}

\begin{pasos}
\paso \celda{C8}: el pago fijo mensual de cada gestor (es igual para todos).
\paso Columna B: los nombres reales de los 21 gestores. La columna D tiene su territorio y la E muestra la región.
\paso Si algún gestor tiene un pago fijo distinto, escríbalo en su celda de la columna F.
\end{pasos}

\section{Borre los ejemplos y compruebe}

\begin{pasos}
\paso Borre los contenedores de ejemplo y las bases de comercialización como en el ejercicio~8 (sección~\ref{sec:ej8}).
\paso Vaya a \hoja{INICIO}: contenedores registrados 0, cuadres OK y conciliación CUADRA.
\paso Guarde. \textbf{El libro está listo para registrar contenedores reales.}
\end{pasos}

\begin{nota}
Para comprobar que todo está bien configurado, registre un contenedor real cuyo resultado conozca aproximadamente y compare. Si el resultado se
aleja mucho de lo esperado, revise en \hoja{CALCULO} qué etapa tiene un costo extraño. Casi siempre es un precio o una norma mal escritos.
\end{nota}

% ============================================================================
\chapter{El trabajo con cada contenedor}\label{cap:diario}
% ============================================================================

\section{El ciclo de un contenedor}

\begin{figure}[H]\centering
\begin{tikzpicture}[
  m/.style={draw=azul,fill=azulclaro,rounded corners=3pt,text width=4.3cm,align=center,minimum height=2.1cm,font=\small},
  f/.style={-{Stealth[length=3mm]},very thick,azul}]
\node[m] (a) {\textbf{Momento 1: llega}\\[2pt]ID, fecha, transitaria, tipo, bultos y \textbf{kg por región}\\Estado: \textbf{Abierto}};
\node[m,right=1.1cm of a] (b) {\textbf{Durante la operación}\\[2pt]Guarde los vales, las hojas de ruta y los comprobantes};
\node[m,right=1.1cm of b] (c) {\textbf{Momento 2: termina}\\[2pt]Datos reales, fecha de fin y documentos\\Estado: \textbf{Cerrado}};
\draw[f] (a) -- (b); \draw[f] (b) -- (c);
\end{tikzpicture}
\caption{Los dos momentos de registro de cada contenedor}
\end{figure}

\section{Momento 1: cuando llega el contenedor}

\begin{pasos}
\paso Abra el archivo maestro y vaya a \hoja{ENTRADA}.
\paso Busque la primera fila vacía: haga clic en \celda{B6} y pulse \tecla{Ctrl}+\tecla{$\downarrow$}. Excel salta al último contenedor
      registrado; la fila siguiente es la libre. Si la hoja está vacía, use la fila 6.
\paso Complete de \textbf{B} a \textbf{I}: ID, fecha de arribo, transitaria, tipo, bultos y kg de Occidente, Centro y Oriente.
\paso En \textbf{Z} ponga \textbf{Abierto}.
\paso Mire las columnas \textbf{AB} y \textbf{AC}: alertas y primera estimación de la utilidad, con las normas.
\paso Guarde.
\end{pasos}

\begin{consejo}
Si una región no recibe carga en ese contenedor, escriba \escriba{0} en sus kg o déjelos vacíos: el libro no le cargará viajes ni combustible a esa región.
\end{consejo}

\section{Momento 2: cuando termina la distribución}

\begin{pasos}
\paso Busque la fila del contenedor. Si hay muchos, use el \textbf{filtro}: en la fila 5, pulse la flecha del título \textbf{Estado}, deje marcado
      solo «Abierto» y pulse \boton{Aceptar}. Para volver a verlos todos, vuelva a la flecha y elija «Seleccionar todo».
\paso Con los documentos delante, escriba los datos reales en las columnas opcionales (J a U), en la estadía (N) y en otros gastos e ingresos (V a X).
\paso Escriba la \textbf{fecha de fin} (Y), cambie el \textbf{Estado} (Z) a \textbf{Cerrado} y anote los documentos (AA).
\paso Revise la alerta (AB). Si aparece, investíguela (sección~\ref{sec:ej3}).
\paso Guarde.
\end{pasos}

\section{Qué escribir en cada columna opcional}

\begin{center}
\begin{tabularx}{\linewidth}{@{}l>{\raggedright\arraybackslash}p{4.2cm}>{\raggedright\arraybackslash}X@{}}
\toprule
\textbf{Col.} & \textbf{Dato} & \textbf{Qué escribir} \\ \midrule
J & Viajes al puerto & Cuántas veces fue el camión propio al puerto por este contenedor (normalmente 1). \\
K & Litros reales al puerto & El \textbf{total} de litros de esos viajes, según los vales. \\
L & Dieta pagada & El \textbf{total} pagado de dieta al chofer, en USD. \\
M & Jornales del desagrupe & Trabajadores $\times$ días. Ejemplo: 4 trabajadores durante 3 días = \escriba{12}. \\
N & Estadía / almacenaje & Lo que cobró la naviera o la transitaria por pasarse de los días libres. \\
O & Litros de distribución en Occidente & Total de litros de la distribución en Occidente. \\
P / S & Litros de traslado a Centro / Oriente & Total de litros del viaje hasta el centro transitorio. \\
Q / T & Litros de distribución en Centro / Oriente & Total de litros de la distribución desde el centro transitorio. \\
R / U & Jornales del centro transitorio & Trabajadores $\times$ días del desagrupe en el centro transitorio de esa región. \\
\bottomrule
\end{tabularx}
\end{center}

\begin{atencion}
Escriba \textbf{totales del contenedor}, no valores por viaje: si hubo 2 viajes de 250 L a Centro, escriba \escriba{500} en P.
Si no conoce un dato real, deje la celda \textbf{vacía}, nunca con 0: un 0 significa «no se gastó nada» y daría una utilidad falsa.
\end{atencion}

\section{Casos especiales}

\begin{description}[style=nextline,leftmargin=1.2em]
  \item[Corregir un contenedor ya cerrado] Corrija la celda equivocada. Todo se recalcula solo. Anote la corrección en Observaciones.
  \item[Un contenedor de otro mes que llega tarde al registro] Regístrelo con su fecha real de arribo. Cambiará el reparto de los gastos fijos de ese mes,
        lo cual es correcto.
  \item[Gastos del contenedor que no tienen columna] Use «Otros gastos directos» (V) y explíquelos en «Detalle» (W).
  \item[Ingresos extra] Domicilios, sobrepeso, etc.: use «Otros ingresos» (X).
  \item[Contenedor que va a una sola región] Ponga todos sus kg en esa región y deje las otras vacías.
\end{description}

% ============================================================================
\chapter{El cierre de cada mes}\label{cap:cierre}
% ============================================================================

El cierre se hace el \textbf{primer día hábil del mes siguiente} y lleva unos 30 minutos. Al terminarlo, los resultados del mes quedan definitivos.
El ejercicio~5 (sección~\ref{sec:ej5}) es un cierre completo de práctica.

\section{Lista de comprobación del cierre}

\begin{center}
\renewcommand{\arraystretch}{1.35}
\begin{tabularx}{\linewidth}{@{}c>{\raggedright\arraybackslash}X l@{}}
\toprule
& \textbf{Tarea} & \textbf{Hoja} \\ \midrule
$\square$ & 1. Todos los contenedores del mes están registrados y en \textbf{Cerrado}. & \hoja{ENTRADA} \\
$\square$ & 2. Las alertas del mes están resueltas o explicadas en Observaciones. & \hoja{ENTRADA} \\
$\square$ & 3. Los gastos fijos reales del mes que difieren del valor base están escritos en la columna del mes. & \hoja{COSTOS\_FIJOS} \\
$\square$ & 4. Si cambió algún precio durante el mes, está como una nueva fila con su fecha. & \hoja{PRECIOS} \\
$\square$ & 5. El mes a liquidar está elegido y la base variable de cada gestor, escrita. & \hoja{COMERCIAL} \\
$\square$ & 6. La conciliación dice \textbf{CUADRA}, o la diferencia está explicada (cobros pendientes). & \hoja{COMERCIAL} \\
$\square$ & 7. Los cuadres dicen \textbf{OK} y la columna «Control» de \hoja{RESUMEN\_MES} es 0. & \hoja{INICIO} \\
$\square$ & 8. Resultado del mes revisado: utilidad, costo por kg, punto de equilibrio, regiones. & \hoja{RESUMEN\_MES} \\
$\square$ & 9. Informe del mes exportado a PDF. & \hoja{DASHBOARD} \\
$\square$ & 10. Archivo guardado y copia de respaldo del mes hecha. & --- \\
\bottomrule
\end{tabularx}
\end{center}
El anexo~\ref{anx:checklist} tiene esta lista en una página aparte para imprimir.

\section{Detalle de cada tarea}

\subsection*{Gastos fijos reales (tarea 3)}
En \hoja{COSTOS\_FIJOS} cada mes copia el valor base. Compare con las facturas y nóminas del mes. Si un gasto fue distinto, escriba el importe real
en la columna de ese mes (los meses empiezan en la columna G, con el primer mes del período).

\begin{atencion}
\textbf{No cambie el «Valor base» (columna F) para corregir un solo mes}: cambiaría todos los meses, incluidos los ya cerrados.
Cambie el valor base solo si el gasto cambia \textbf{desde ahora y para siempre}, y antes siga el procedimiento del capítulo~\ref{cap:cambios}.
\end{atencion}

\subsection*{Liquidación de los gestores (tareas 5 y 6)}
\begin{pasos}
\paso En \hoja{COMERCIAL}, escriba en \celda{C4} el primer día del mes que cierra.
\paso En la columna \textbf{G} escriba lo que gestionó cada gestor en el mes: USD cobrados o kg, según la base.
\paso La columna \textbf{I} es lo que hay que pagar a cada uno y la fila 33, el total.
\paso Mire la conciliación (filas 36 a 39). Si dice REVISAR, compare gestor por gestor con los contenedores del mes.
      Las causas típicas son un cobro sin declarar, un contenedor sin registrar o un cobro pendiente.
\end{pasos}

\begin{nota}
Si la empresa paga la comisión sobre lo \textbf{efectivamente cobrado} (es lo recomendable), puede que al cierre falte cobrar algo.
En ese caso la diferencia de la conciliación es lo pendiente de cobro: anótelo y compruebe que se cobre el mes siguiente.
\end{nota}

\section{La copia de respaldo}\label{sec:respaldo}

Como usted es la única persona que lleva el libro, la copia de respaldo es su seguro: si el archivo se daña, se borra o se rompe la computadora,
el trabajo no se pierde.

\begin{pasos}
\paso Guarde el archivo maestro (\tecla{Ctrl}+\tecla{G}) y ciérrelo.
\paso Abra la carpeta del archivo en el Explorador de Windows. Haga clic derecho sobre el archivo $\rightarrow$ \menu{Copiar}, y luego clic derecho en
      un espacio vacío $\rightarrow$ \menu{Pegar}.
\paso Cambie el nombre de la copia (clic derecho $\rightarrow$ \menu{Cambiar nombre}) por \escribal{Respaldo\_\allowbreak 2026-10\_\allowbreak Rentabilidad.xlsx}
      (año y mes del cierre).
\paso Guarde también una copia \textbf{fuera de la computadora}: memoria USB, correo electrónico o la nube.
\end{pasos}

\begin{consejo}
Haga además una copia rápida cada viernes. Si un día comete un error que no puede deshacer, recupere la última copia y vuelva a escribir solo lo de
los últimos días. \textbf{Nunca trabaje sobre una copia de respaldo}: siempre sobre el maestro.
\end{consejo}

% ============================================================================
\chapter{Cuando algo cambia}\label{cap:cambios}
% ============================================================================

La regla general: \textbf{lo que cambia desde una fecha no debe alterar el pasado}. Cada caso tiene su forma correcta.

\section{Cambia un precio o una tarifa}
Diésel, tarifa por kg, cita, dieta, pago por jornal, servicios de transitaria, comisión\dots
\begin{pasos}
\paso En \hoja{PRECIOS}, copie la última fila con precios en la fila siguiente vacía.
\paso Cambie la fecha de la columna A por la fecha desde la que rige el cambio.
\paso Cambie solo el precio que varió y anote el documento en la columna O.
\end{pasos}
Es el ejercicio~4 (sección~\ref{sec:ej4}). Mantenga las fechas \textbf{en orden}, de la más antigua a la más reciente.

\section{Cambia un gasto fijo de forma permanente}
Por ejemplo, sube el salario de la oficina desde marzo.
\begin{pasos}
\paso En \hoja{COSTOS\_FIJOS}, vaya a la fila del concepto y a la columna de \textbf{marzo}.
\paso Escriba el nuevo importe. Cópielo con \tecla{Ctrl}+\tecla{C}, seleccione las celdas de los meses siguientes y péguelo con \tecla{Ctrl}+\tecla{V}.
\paso \textbf{No toque el valor base} (columna F): así los meses anteriores conservan el importe antiguo.
\end{pasos}

\section{Cambia una norma técnica}
Si después de varios contenedores comprueba que una norma ya no es real (por ejemplo, el camión consume 70 L y no 60), cámbiela en \hoja{PARAMETROS}.
La norma solo afecta a las celdas opcionales que estén \textbf{vacías}. Por eso conviene escribir siempre los litros reales: así un cambio de norma
nunca altera el pasado.

\section{Se compra, se vende o se retira un vehículo o equipo}
\begin{itemize}
  \item \textbf{Compra}: agréguelo en una fila libre de \hoja{ACTIVOS} con su fecha de alta. La depreciación empieza ese mes.
  \item \textbf{Venta o retiro}: escriba la \textbf{fecha de baja} en la columna K. Desde el mes siguiente deja de depreciarse y de cargar mantenimiento.
        \textbf{No borre la fila}: los meses anteriores la necesitan.
\end{itemize}

\section{Cambia un gestor}
\begin{itemize}
  \item \textbf{Sustitución} en el mismo territorio: cambie el nombre en la columna B de \hoja{COMERCIAL}.
  \item \textbf{Cambio de pago fijo para todos}: para que no cambien los meses anteriores, siga la sección «Cambia un gasto fijo de forma
        permanente»: escriba el nuevo total de la fila 8 de \hoja{COSTOS\_FIJOS} en los meses desde el cambio. Luego actualice \celda{C8} de \hoja{COMERCIAL}.
\end{itemize}

\section{Cambian los impuestos}
Actualice las tasas en \hoja{PARAMETROS} (sección D) \textbf{después de confirmarlo con el contador}.

\begin{atencion}
Las tasas de \hoja{PARAMETROS} se aplican a \textbf{todos} los meses del libro. Si cambian a mitad de año, guarde antes una copia del libro con
las tasas antiguas como histórico, o anote el cambio para tenerlo en cuenta en la liquidación anual.
\end{atencion}

\section{Se llena el libro: nuevo período}
El libro admite \textbf{24 meses} y \textbf{300 contenedores}. Al acercarse a cualquiera de los dos límites:
\begin{pasos}
\paso Haga un respaldo y guarde el archivo como histórico, por ejemplo \escribal{Rentabilidad\_\allowbreak HISTORICO\_\allowbreak 2026-2027.xlsx}.
\paso Haga una copia como nuevo maestro.
\paso En el nuevo maestro, cambie \texttt{Fecha\_Inicio} en \hoja{PARAMETROS} al primer mes del nuevo período.
\paso Borre los contenedores de \hoja{ENTRADA} (con \tecla{Supr}) y revise que \hoja{PRECIOS}, \hoja{ACTIVOS} y \hoja{COSTOS\_FIJOS}
      tengan los valores vigentes.
\end{pasos}

\section{Ver o cambiar algo protegido}
Las hojas están protegidas \textbf{sin contraseña}, solo para evitar errores. Si necesita cambiar algo bloqueado:
\menu{Revisar $\rightarrow$ Desproteger hoja}. Haga el cambio y vuelva a protegerla con \menu{Revisar $\rightarrow$ Proteger hoja}
$\rightarrow$ \boton{Aceptar}, sin escribir contraseña.

\begin{atencion}
Con la hoja desprotegida, las fórmulas se pueden borrar sin querer. Si algo falla después de un cambio así, recupere la última copia de respaldo.
\end{atencion}

% ============================================================================
\chapter{Entender los resultados}\label{cap:resultados}
% ============================================================================

Este capítulo explica, con los números de ejemplo del libro, qué significa cada indicador y qué decisiones ayuda a tomar.

\section{De los ingresos a la utilidad}

\begin{figure}[H]\centering
\begin{tikzpicture}[
  b/.style={draw=azul,minimum height=0.9cm,text width=6.3cm,align=left,font=\small,inner sep=4pt},
  r/.style={b,fill=azulclaro,font=\small\bfseries}]
\node[r] (i) {Ingresos = kg $\times$ tarifa de cada región};
\node[b,below=0pt of i] (v) {$-$ Costos variables (puerto, transitaria, distribución, comisión variable, reclamaciones, banca)};
\node[r,below=0pt of v] (m) {= Margen de contribución};
\node[b,below=0pt of m] (f) {$-$ Costos fijos del mes (o la parte de cada contenedor)};
\node[b,below=0pt of f] (t) {$-$ Tributos sobre ingresos (10\,\% + 1\,\%)};
\node[r,below=0pt of t] (u) {= Utilidad antes de impuesto (UAI)};
\node[b,below=0pt of u] (iu) {$-$ Impuesto sobre utilidades (35\,\%, estimado)};
\node[r,below=0pt of iu,fill=green!15] (n) {= Utilidad neta estimada};
\node[right=0.6cm of m,text width=6cm,font=\small,align=left] {Con los datos de ejemplo:\\ingresos \vIngtotal; margen de contribución \vMCpct{} de los ingresos.};
\node[right=0.6cm of u,text width=6cm,font=\small,align=left] {UAI \vUAItotal};
\node[right=0.6cm of n,text width=6cm,font=\small,align=left] {Utilidad neta \vUNtotal};
\end{tikzpicture}
\caption{Cómo se forma la utilidad}
\end{figure}

\section{Los indicadores, uno por uno}

\begin{description}[style=nextline,leftmargin=1.2em,itemsep=6pt]
\item[Margen de contribución (MC)]
  Lo que queda de los ingresos después de pagar los gastos que dependen de cada contenedor. Es el dinero disponible para pagar los gastos fijos y
  para ganar. Un MC del \vMCpct{} quiere decir que ese porcentaje de cada dólar que entra queda disponible para cubrir los gastos fijos y dar ganancia.
\item[Utilidad antes de impuesto (UAI)]
  La ganancia del período después de todos los gastos y de los tributos sobre los ingresos. Si es negativa (en rojo y entre paréntesis), el
  período tuvo pérdidas.
\item[Utilidad neta estimada]
  La UAI menos el impuesto sobre utilidades estimado (35\,\%). El impuesto real se liquida \textbf{una vez al año}: use la tabla anual de \hoja{RESUMEN\_MES}
  y confírmela con el contador.
\item[Ingreso por kg y costo total por kg]
  Cuánto entra y cuánto cuesta, en promedio, cada kg transportado. En el ejemplo entran \vIngKg{} y cuesta \vCostoKg{} por kg. La diferencia es la
  utilidad por kg. Si el costo por kg sube de un mes a otro, busque la causa en la estructura de costos del \hoja{DASHBOARD}.
\item[Punto de equilibrio]
  Los kg (o contenedores) necesarios para cubrir \textbf{exactamente} todos los gastos fijos: por debajo hay pérdidas; por encima, ganancias.
  En el ejemplo hacen falta \vPEkg{} kg en el período.
\item[Margen de seguridad]
  Cuánto pueden bajar los kg antes de entrar en pérdidas. Un \vMSeg{} significa que el volumen podría caer ese porcentaje y la empresa todavía
  no perdería. Cuanto mayor, más tranquilidad.
\item[Fijos no absorbidos]
  Gastos fijos que ningún contenedor cubrió: un mes sin arribos, o una región sin carga ese mes. Indican \textbf{capacidad ociosa}, es decir,
  se pagó una estructura que no se usó.
\item[Desvío del combustible]
  Cuánto se apartó el consumo real de la norma. Positivo: se gastó más de lo normal.
\end{description}

\section{Leer la rentabilidad por región}

La utilidad de cada región se calcula así: sus ingresos, menos sus propios costos (combustible, desagrupe regional, comisión y gastos fijos de la
región), menos \textbf{su parte de los costos compartidos} (puerto, transitaria y gastos fijos generales) según sus kg.

\begin{ejemplo}
Con los datos de ejemplo, Occidente gana \vUAIocc, Centro \vUAIcen{} y Oriente \textcolor{red}{\vUAIori}. Oriente cobra lo mismo por kg, pero tiene
dos tramos de transporte y un segundo desagrupe. Con una tarifa única, \textbf{Occidente está pagando las pérdidas de Oriente}.
El \hoja{SIMULADOR} calcula la tarifa mínima de Oriente (\vSimTminOri{} por kg) y el ejercicio~7 muestra el efecto de subirla.
\end{ejemplo}

\section{Preguntas que el libro le ayuda a responder}

\begin{center}
\begin{tabularx}{\linewidth}{@{}>{\raggedright\arraybackslash}p{6.2cm}>{\raggedright\arraybackslash}X@{}}
\toprule
\textbf{Pregunta} & \textbf{Dónde mirar} \\ \midrule
¿Ganamos dinero este mes? & \hoja{RESUMEN\_MES}, columna «Utilidad antes de impuesto» \\
¿Cuál fue el contenedor menos rentable y por qué? & \hoja{ENTRADA}, columna AC; detalle en \hoja{CALCULO} \\
¿Qué región conviene revisar? & \hoja{DASHBOARD}, rentabilidad por región \\
¿Qué tarifa mínima debemos cobrar? & \hoja{SIMULADOR}, fila «Tarifa mínima de equilibrio» \\
¿Cuántos contenedores necesitamos al mes? & \hoja{SIMULADOR}, punto de equilibrio y tabla de contenedores \\
¿Qué pasa si sube el diésel? & \hoja{SIMULADOR}, tabla de sensibilidad \\
¿En qué se va el dinero? & \hoja{DASHBOARD}, estructura de costos \\
¿Cuánto pagaremos de impuesto este año? & \hoja{RESUMEN\_MES}, liquidación anual estimada \\
\bottomrule
\end{tabularx}
\end{center}

\section{Cómo presentar los resultados a la dirección}
\begin{pasos}
\paso En \hoja{DASHBOARD}, elija el período.
\paso Exporte la hoja a PDF: \menu{Archivo $\rightarrow$ Exportar $\rightarrow$ Crear PDF}, o \menu{Imprimir} y elija la impresora «PDF».
\paso Acompáñelo con tres frases: la utilidad del período y su tendencia, la región más débil y su causa, y una propuesta (tarifa, consolidar viajes, controlar combustible\dots).
\end{pasos}

% ============================================================================
\chapter{Solución de problemas}\label{cap:problemas}
% ============================================================================

\section{Mensajes y situaciones frecuentes}

\renewcommand{\arraystretch}{1.3}
\begin{longtable}{@{}>{\raggedright\arraybackslash}p{4.4cm}>{\raggedright\arraybackslash}p{4.6cm}>{\raggedright\arraybackslash}p{6.4cm}@{}}
\toprule
\textbf{Lo que ve} & \textbf{Por qué ocurre} & \textbf{Qué hacer} \\ \midrule
\endhead
«La celda o el gráfico que intenta cambiar está en una hoja protegida» & Intentó escribir en una celda de fórmula. & \boton{Aceptar}. Escriba solo en las celdas amarillas. \\
«Número no válido», «Fecha no válida» o «Valor no válido» & Escribió texto donde va un número, una fecha imposible o un valor que no está en la lista. & \boton{Reintentar} y escriba solo el número, la fecha dd/mm/aaaa o un valor de la lista. \\
No puede escribir en nada; hay una barra amarilla arriba & El archivo está en «Vista protegida». & Pulse \boton{Habilitar edición}. \\
\texttt{\#\#\#\#\#\#} en una celda & La columna es demasiado estrecha para el número. & Haga doble clic en el borde derecho de la letra de la columna, o reduzca el zoom. \\
Una fecha queda alineada a la izquierda & Excel la tomó como texto: formato o separador incorrecto. & Escríbala como \escriba{27/10/2026}. Si sigue igual, pruebe con el formato de fecha de su Windows. \\
Un número queda alineado a la izquierda & Se tomó como texto: tiene letras, espacios o el separador decimal equivocado. & Vuelva a escribirlo solo con cifras. Pruebe con punto o con coma decimal. \\
Una fila de \hoja{ENTRADA} no muestra resultado & Falta el ID en la columna B. & Escriba el ID del contenedor. \\
Alerta «Falta fecha» o «Falta transitaria» & Dato obligatorio vacío. & Complete la columna C o D. \\
Alerta «Mes fuera del rango de COSTOS\_FIJOS» & La fecha cae antes de Fecha\_Inicio o después de los 24 meses. & Revise la fecha. Si es correcta, vea «Se llena el libro» (capítulo~\ref{cap:cambios}). \\
Alerta «Fecha anterior a la 1ª vigencia de precios» & No hay precios para esa fecha. & En \hoja{PRECIOS}, ponga en \celda{A6} una fecha más antigua. \\
Alerta «Combustible \dots{} sobre norma» & El consumo real supera la norma más la tolerancia. & Investigue la causa (sección~\ref{sec:ej3}). Corrija el dato o explíquelo en Observaciones. \\
Alerta «Contenedor con pérdida» & Los costos del contenedor superan sus ingresos. & Revise en \hoja{CALCULO} qué etapa pesa más. Puede deberse a pocos contenedores en el mes. \\
La utilidad de un contenedor cambió sin tocarlo & Se registró otro contenedor del mismo mes y los gastos fijos se repartieron de nuevo. & Es normal (ejercicio~2). Es definitiva al cerrar el mes. \\
Conciliación: «REVISAR» & Lo declarado por los gestores no coincide con los contenedores del mes. & Busque el cobro sin declarar o el contenedor sin registrar (ejercicio~5). \\
Cuadres de \hoja{INICIO}: «REVISAR» & Algo se dañó, normalmente una fórmula borrada con la hoja desprotegida. & \tecla{Ctrl}+\tecla{Z}. Si no basta, recupere la última copia de respaldo. \\
Borré sin querer el contenido de una celda de un mes en \hoja{COSTOS\_FIJOS} & La celda copiaba el valor base. & Escriba \escriba{=\$F10}, cambiando 10 por el número de la fila, o escriba el importe del mes. \\
Eliminé una fila o columna (con la hoja desprotegida) & Las fórmulas pierden sus referencias (\texttt{\#¡REF!}). & \tecla{Ctrl}+\tecla{Z} enseguida. Si ya guardó, recupere la última copia de respaldo. \\
\texttt{\#¡VALOR!}, \texttt{\#¡DIV/0!} o \texttt{\#N/A} & Un dato tiene un formato inválido o falta una referencia. & Revise el último dato que escribió. Si no lo encuentra, recupere la copia de respaldo. \\
El archivo no abre o está dañado & Corte de corriente o fallo del disco. & Abra la última copia de respaldo y vuelva a escribir lo posterior. \\
En LibreOffice algo se ve distinto & Diferencias de presentación entre programas. & Los cálculos son los mismos. Si un gráfico se ve raro, ábralo en Excel. \\
\bottomrule
\end{longtable}

\section{Si nada de lo anterior funciona}
\begin{pasos}
\paso No guarde. Cierre el archivo sin guardar los cambios.
\paso Ábralo de nuevo: tendrá la última versión guardada.
\paso Si el problema sigue, abra la copia de respaldo más reciente, guárdela como nuevo maestro y repita el trabajo posterior a esa copia.
\end{pasos}

\appendix
% ============================================================================
\chapter{Columnas de la hoja ENTRADA}\label{anx:columnas}
% ============================================================================

\renewcommand{\arraystretch}{1.25}
\begin{longtable}{@{}l>{\raggedright\arraybackslash}p{3.9cm}>{\raggedright\arraybackslash}p{4.9cm}>{\raggedright\arraybackslash}p{3.6cm}c@{}}
\toprule
\textbf{Col.} & \textbf{Nombre} & \textbf{Qué escribir} & \textbf{Documento fuente} & \textbf{Oblig.} \\ \midrule
\endhead
A & Nº & Nada: se numera solo & --- & --- \\
B & ID contenedor & Número del contenedor o del BL & BL / manifiesto & Sí \\
C & Fecha de arribo & dd/mm/aaaa & Aviso de arribo & Sí \\
D & Transitaria & ENCI, A.V u Otra (lista) & Documentos de la transitaria & Sí \\
E & Tipo & 20', 40', 40' HC u Otro (lista) & BL & No \\
F & Nº de bultos & Cantidad de bultos & Manifiesto & No \\
G & kg Occidente & kg con destino a Occidente & Desglose por destinatario & Sí \\
H & kg Centro & kg con destino a Centro & Desglose por destinatario & Sí \\
I & kg Oriente & kg con destino a Oriente & Desglose por destinatario & Sí \\
J & Viajes al puerto & Nº de viajes del camión propio & Hoja de ruta & Opc. \\
K & Litros reales al puerto & Total de litros & Vales de combustible & Opc. \\
L & Dieta pagada & Total en USD & Comprobante de dieta & Opc. \\
M & Jornales del desagrupe & Trabajadores $\times$ días & Nómina / control de asistencia & Opc. \\
N & Estadía / almacenaje & USD cobrados por demora & Factura de la naviera o de la transitaria & No \\
O & Litros distribución Occidente & Total de litros & Vales / hoja de ruta & Opc. \\
P & Litros traslado a Centro & Total de litros & Vales / hoja de ruta & Opc. \\
Q & Litros distribución Centro & Total de litros & Vales / hoja de ruta & Opc. \\
R & Jornales centro transitorio Centro & Trabajadores $\times$ días & Nómina / asistencia & Opc. \\
S & Litros traslado a Oriente & Total de litros & Vales / hoja de ruta & Opc. \\
T & Litros distribución Oriente & Total de litros & Vales / hoja de ruta & Opc. \\
U & Jornales centro transitorio Oriente & Trabajadores $\times$ días & Nómina / asistencia & Opc. \\
V & Otros gastos directos & USD & Factura / comprobante & No \\
W & Detalle de otros gastos & Texto breve & --- & No \\
X & Otros ingresos & USD & Factura emitida & No \\
Y & Fecha de fin de la distribución & dd/mm/aaaa & Última entrega & No \\
Z & Estado & Abierto o Cerrado (lista) & --- & Sí \\
AA & Observaciones & Números de documentos, explicaciones & --- & No \\
AB & Alertas & Nada: automático & --- & --- \\
AC & Utilidad antes de impuesto & Nada: automático & --- & --- \\
\bottomrule
\end{longtable}
\noindent\textbf{Opc.} = opcional: si se deja vacía, se usa la norma técnica de \hoja{PARAMETROS}.

% ============================================================================
\chapter{Glosario}\label{anx:glosario}
% ============================================================================

\begin{description}[style=nextline,leftmargin=1.2em,itemsep=4pt]
\item[Absorción (costeo por)] Repartir los gastos fijos entre los contenedores para conocer su costo completo.
\item[Alerta] Aviso automático de que falta un dato o de que algo se sale de lo normal.
\item[Celda] Cada casilla de una hoja, nombrada por su columna y su fila, como \celda{C5}.
\item[Centro de costo] Grupo al que se asigna un gasto: GENERAL, Occidente, Centro u Oriente.
\item[Conciliación] Comparación de dos registros que deberían coincidir: lo que declaran los gestores y lo que dicen las operaciones.
\item[Costo fijo] Gasto que se paga igual haya o no contenedores: alquiler, salarios de oficina, depreciación\dots
\item[Costo variable] Gasto que depende de cada contenedor: combustible, jornales, cita, comisión variable\dots
\item[Depreciación] Pérdida de valor del camión o de un equipo por el uso y el tiempo, repartida por meses.
\item[Desagrupe] Separar la carga del contenedor por destinatario.
\item[Fijos no absorbidos] Gastos fijos de un mes que ningún contenedor cubrió.
\item[Jornal] Un día de trabajo de una persona. 4 personas durante 3 días son 12 jornales.
\item[Margen de contribución] Ingresos menos costos variables.
\item[Margen de seguridad] Cuánto pueden bajar las ventas antes de tener pérdidas.
\item[Norma técnica] Consumo considerado normal (litros por viaje, kg por jornal\dots). Se usa cuando falta el dato real y para generar alertas.
\item[Período] Los 24 meses que abarca el libro, desde \texttt{Fecha\_Inicio}.
\item[Prorrateo] Reparto proporcional de un gasto, por ejemplo según los kg.
\item[Punto de equilibrio] Volumen con el que los ingresos cubren exactamente todos los gastos: ni se gana ni se pierde.
\item[Tarifa mínima de equilibrio] Precio por kg más bajo con el que una región no pierde dinero.
\item[UAI] Utilidad antes del impuesto sobre utilidades.
\item[Vigencia] Período durante el que rige un precio. Empieza en la fecha de su fila en \hoja{PRECIOS}.
\end{description}

% ============================================================================
\chapter{Lista de comprobación mensual (para imprimir)}\label{anx:checklist}
% ============================================================================

\noindent Mes que se cierra: \rule{5cm}{0.4pt} \hfill Fecha del cierre: \rule{4cm}{0.4pt}

\bigskip
\renewcommand{\arraystretch}{2.0}
\begin{tabularx}{0.96\linewidth}{@{}c>{\raggedright\arraybackslash}X p{3.4cm}@{}}
\toprule
& \textbf{Tarea} & \textbf{Observaciones} \\ \midrule
$\square$ & Todos los contenedores del mes registrados y en «Cerrado» (\hoja{ENTRADA}) & \\
$\square$ & Alertas resueltas o explicadas & \\
$\square$ & Gastos fijos reales del mes escritos (\hoja{COSTOS\_FIJOS}) & \\
$\square$ & Cambios de precio registrados como filas nuevas (\hoja{PRECIOS}) & \\
$\square$ & Base variable de los 21 gestores escrita (\hoja{COMERCIAL}) & \\
$\square$ & Conciliación de comercialización: CUADRA o diferencia explicada & \\
$\square$ & Cuadres de \hoja{INICIO}: OK y OK & \\
$\square$ & Resultado revisado: utilidad, costo/kg, equilibrio, regiones & \\
$\square$ & Informe del \hoja{DASHBOARD} exportado a PDF & \\
$\square$ & Archivo guardado y copia de respaldo hecha (USB, correo o nube) & \\
\bottomrule
\end{tabularx}

\vspace{1cm}
\noindent Utilidad antes de impuesto del mes: \rule{4cm}{0.4pt}\hfill Contenedores del mes: \rule{3cm}{0.4pt}

\vspace{14pt}
\noindent Firma: \rule{5cm}{0.4pt}

% ============================================================================
\chapter{Cómo actualizar esta guía}\label{anx:mantener}
% ============================================================================

Esta guía está escrita en LaTeX para que cualquier cambio sea sencillo y la guía siga siempre coherente con el libro de Excel.

\section{Dónde está cada cosa}
\begin{center}
\begin{tabularx}{\linewidth}{@{}l X@{}}
\toprule
\textbf{Archivo} & \textbf{Contenido} \\ \midrule
\texttt{guia/guia\_usuario.tex} & Archivo principal: portada, orden de los capítulos y versión (\verb|\version|). \\
\texttt{guia/capitulos/*.tex} & Un archivo por capítulo. Para cambiar un texto, edite el capítulo correspondiente. \\
\texttt{guia/estilo.sty} & Colores, cuadros (nota, atención, consejo, resultado) y comandos (\verb|\hoja|, \verb|\celda|, \verb|\escriba|\dots). \\
\texttt{guia/valores.tex} & Números de los ejercicios. \textbf{Se genera solo}: no lo edite a mano. \\
\texttt{guia/img/} & Capturas del libro. \textbf{Se generan solas}. \\
\texttt{build/capturas.py} & Define los ejercicios, las capturas (hoja, rango, recuadros) y los valores citados. \\
\bottomrule
\end{tabularx}
\end{center}

\section{Volver a generar la guía}
Requiere Python con \texttt{openpyxl}, \texttt{lxml} y \texttt{Pillow}, LibreOffice Calc, \texttt{pdftoppm} y una distribución de LaTeX (TeX Live o MiKTeX).
\begin{verbatim}
build/construir.sh          # 1. regenera el libro de Excel
build/construir_guia.sh     # 2. capturas + valores + PDF
\end{verbatim}
Si cambia el libro de Excel (una hoja, una columna, una fórmula), ejecute ambos pasos: las imágenes y los números de la guía se actualizan solos.
Si un cambio mueve una hoja o una columna, ajuste también los rangos de \texttt{build/capturas.py} y las celdas que nombra el texto.

\section{Comandos de la guía}
\begin{center}
\begin{tabular}{@{}p{7.2cm}p{8.4cm}@{}}
\toprule
\textbf{Escriba en el .tex} & \textbf{Resultado} \\ \midrule
\verb|\hoja{ENTRADA}| & \hoja{ENTRADA} \\
\verb|\celda{C5}| & \celda{C5} \\
\verb|\escriba{5000}| & \escriba{5000} \\
\verb|\tecla{Enter}| & \tecla{Enter} \\
\verb|\recuadro{1}| & \recuadro{1} \\
\verb|\captura{ej2_entrada}{Pie}| & Inserta \texttt{img/ej2\_entrada.png} con su pie de figura \\
\verb|\begin{resultado}...\end{resultado}| & Cuadro verde «Lo que debe ver si lo hizo bien» \\
\bottomrule
\end{tabular}
\end{center}

\end{document}
